% Opuscules de François Maurolyco — Latin 17859, lecture modernisée du volume entier.
%
% Pass: Opus 5 (claude-opus-5), 15 September 2026 — first pass, unchecked
% against the views by a human, and an interpretation rather than a record.
% Where this file and the transcriptions differ in substance, the
% transcriptions (batch-01.fr.tex … batch-03.fr.tex) are the record.
%
% Sources read: the three batch transcriptions of this volume and nothing
% else. No image was consulted, and no printed edition of Maurolico — neither
% the Messina edition of 1558 from which the copy was made, nor the Opuscula
% mathematica of 1575, nor any modern edition — nor any other transcription.
% Where the transcription has a gap, this reading says so and does not fill
% it. All three batches were transcribed on Opus 5; together they cover
% views 1–49, the whole volume (view 1, the board and the certificate, is
% not transcribed; view 49 right is blank).
%
% What the volume turned out to be. A seventeenth-century copy, in one hand,
% of most of the Messina volume of 1558 on the sphere: dedications and the
% index of Maurolico's works; the « De Sphæra Sermo »; Autolycus, Theodosius
% « De Habitationibus » and Euclid's « Phænomena » in Maurolico's versions;
% the tables of sines, tangents (« tabula foecunda ») and secants (« tabula
% benefica ») with eight worked problems; the end of a first book and the
% whole of a second book of Maurolico's own spherical trigonometry; a text on
% Regiomontanus's rules dated August 1530; a dedication dated 1 June 1557;
% and the copyist's note that the « Compendium Mathematicæ » is left out
% because Mersenne had printed it. Theodosius's and Menelaus's Sphaerica are
% announced and not copied.
%
% The order followed. Sections go in the order of the binding, except where
% the transcriptions' own notes show leaves bound out of the text's order:
% the Autolycus page of view 48 L is read before view 10 R; the leaves inked
% 16 and 17 are read in the order of the problems (views 18 R, 19 L, 17 R,
% 18 L); the leaves inked 41–43 (views 42 R – 45 L) are read after view
% 47 L and before view 47 R. One grouping is this reading's own and is said
% to be so: problems « 7 » and « 8 » of views 48 R – 49 L are read after
% problem 6 of view 18 L.
%
% Notation translated here, each with a footnote at its first use: « sinus
% rectus » / « sinus secundus » / « sinus versus » into sin, cos, vers with
% the radius 100 000 taken as 1; « umbra versa » and « umbra recta » with a
% gnomon equal to the radius into tan and cot; the « Radius » of the tabula
% benefica into sec; « arcus circulorum majorum » into arcs of great circles;
% « latitudo ortiva » into the rising amplitude; « differentia ascensionalis »
% into the ascensional difference; « differentia transitus per coeli medium »
% into the difference in right ascension; the composition of ratios into
% products of quotients; the leaf's lower-case figure letters into capitals.
% Modern names — Menelaus's theorem, the sine rule, the spherical law of
% cosines, Napier's rules, the addition formula, Ptolemy's theorem — are
% attached only where the statement coincides, and never as a claim about
% who found what first.
%
% Every worked number and every table entry the transcriptions flag was
% recomputed; the recomputations are this reading's, and each correction
% says in its footnote what the leaf has.
\documentclass[11pt,a4paper]{article}
% The preamble shared by every transcription of the mathematicians' manuscripts in Gallica.
%
% It rests on one idea: **uncertainty compounds, it does not resolve.** A
% transcription that smooths over a doubtful word has destroyed the only thing
% it was made for — a reader must be able to tell, at every point, what was on
% the page and what was supplied. The apparatus macros below are therefore the
% critical apparatus, not decoration.
%
% The same commands are recognised by `scripts/render.mjs`, which produces the
% site's reading view, and by `scripts/tei.mjs`. A macro added here must be
% added there too, or rendering fails loudly — which is the wanted behaviour.
%
% Comments here are English, like the rest of the repository. The documents
% this preamble sets are French, and that is deliberate: see the skills.

\ProvidesPackage{germain}[2026/09/15 Transcriptions of mathematicians' manuscripts in Gallica]

\usepackage{amsmath,amssymb,amsthm}
\usepackage{mathrsfs}
% Kept from the parent preamble so the renderer's diagram support stays
% available; these manuscripts draw no commutative diagrams, and nothing here uses it.
\usepackage{tikz-cd}
\usepackage[english,french]{babel}
\usepackage{fontspec}

% --- Greek ------------------------------------------------------------------
% The text face stays fontspec's default, Latin Modern, which has no Greek:
% XeTeX drops every glyph it lacks with a line in the log and nothing on the
% page, and the Greek words the manuscripts quote vanished from the PDFs. So
% the Greek and Greek Extended blocks get a XeTeX character class of their own,
% and entering or leaving a run of it switches the family to CMU Serif — the
% same Computer Modern design, with polytonic Greek — and back. Only the
% family changes: a Greek word in \textit or \textbf keeps its shape and series.
% The transitions to and from every other class carry the tokens that class
% has towards an ordinary letter, so French spacing before « ; » or inside
% « » still holds after a Greek word. No grouping, so no brace or \endgroup
% can come out unbalanced; maths is left alone, since XeTeX inserts nothing
% there. Kept identical in `exercices.sty`.
\makeatletter
\newfontfamily\germainGreekfont{cmun}[
  Extension=.otf, NFSSFamily=germaingreek,
  UprightFont=*rm, ItalicFont=*ti, SlantedFont=*sl,
  BoldFont=*bx, BoldItalicFont=*bi, BoldSlantedFont=*bl]
\newXeTeXintercharclass\germain@greek
\def\germain@greekrange#1#2{%
  \@tempcnta=#1\relax
  \loop
    \XeTeXcharclass\@tempcnta=\germain@greek
    \ifnum\@tempcnta<#2\advance\@tempcnta\@ne\repeat}
\germain@greekrange{"0370}{"03FF}
\germain@greekrange{"1F00}{"1FFF}
\def\germain@greekprev{\rmdefault}
\def\germain@greekin{%
  \edef\germain@greekprev{\f@family}\fontfamily{germaingreek}\selectfont}
\def\germain@greekout{\fontfamily{\germain@greekprev}\selectfont}
\def\germain@greekjoin{%
  \XeTeXinterchartokenstate=\@ne
  \@tempcnta=\z@
  \loop
    \ifnum\@tempcnta=\germain@greek\else
      \XeTeXinterchartoks\germain@greek\@tempcnta=\expandafter{%
        \expandafter\germain@greekout\the\XeTeXinterchartoks\z@\@tempcnta}%
      \XeTeXinterchartoks\@tempcnta\germain@greek=\expandafter{%
        \the\XeTeXinterchartoks\@tempcnta\z@\germain@greekin}%
    \fi
    \ifnum\@tempcnta<4095\advance\@tempcnta\@ne\repeat}
% babel-french sets its own classes, and saves and restores the interchar
% state, each time a language is selected: join again after it does.
\addto\extrasfrench{\germain@greekjoin}
\addto\extrasenglish{\germain@greekjoin}
\AtBeginDocument{\germain@greekjoin}
\makeatother

\usepackage{geometry}
\geometry{a4paper,margin=25mm,marginparwidth=28mm}
\usepackage{marginnote}
\usepackage{xcolor}
\usepackage[normalem]{ulem}
\setlength{\emergencystretch}{3em}
\usepackage{hyperref}
\hypersetup{colorlinks=true,linkcolor=black,urlcolor=black,pdfborder={0 0 0}}

% --- Batch metadata ---------------------------------------------------------
% Taken as they stand by the rendering script, which shows them at the head of
% the reading view. They fix what the file claims to cover. The macro names are
% the parent project's, so that the scripts stay shared; \folder here means the
% volume's slug — `fr-9115` — and \pages counts Gallica views.

\newcommand{\folder}[1]{\gdef\grothFolder{#1}}
\newcommand{\batch}[1]{\gdef\grothBatch{#1}}
\newcommand{\pages}[2]{\gdef\grothFirst{#1}\gdef\grothLast{#2}}
\newcommand{\foldertitle}[1]{\gdef\grothFoldertitle{#1}}

% The holder's own shelfmark, printed as they write it: « Français 9115 ».
\newcommand{\shelfmark}[1]{\gdef\grothShelfmark{#1}}

% The dating, copied verbatim from the holder's catalogue — never inferred
% here. « XIXe siècle » is the BnF's; a pass that dates a leaf from its content
% says so in a \note{}, as its own inference.
\newcommand{\dating}[1]{\gdef\grothDating{#1}}

% The status notice, printed in the title block and repeated as a diagonal
% watermark on every page of the PDF. The manuscripts are in the public
% domain; what this line declares is not a right but a quality — an
% unverified first pass by a machine. The skills write the line; a file
% without it gets no watermark, so leaving it out is visible in the output.
\newcommand{\watermark}[1]{\gdef\grothWatermark{#1}}

\gdef\grothFolder{?}\gdef\grothBatch{}\gdef\grothFirst{?}\gdef\grothLast{?}%
\gdef\grothFoldertitle{}\gdef\grothShelfmark{}\gdef\grothDating{}\gdef\grothWatermark{}

% --- The résumé -------------------------------------------------------------
\newenvironment{resume}
  {\small\setlength{\parskip}{0.45em}}
  {\par\medskip\hrule\medskip}

% --- Keywords ---------------------------------------------------------------
% In English, closing the modernised reading's résumé: the modern vocabulary
% under which to search for what the leaves construct. The manifest extracts
% them to tag the volume — this line is the single source of the tags.
\newcommand{\keywords}[1]{\par\smallskip\noindent
  {\footnotesize\textsc{Keywords} --- \textit{#1}}\par}

% --- The critical apparatus -------------------------------------------------

% The Gallica view number, in the margin. This is the anchor that lets a
% disagreement over a formula be settled by looking at *one* image rather than
% a volume of 750. The site uses it to turn the facsimile as the transcription
% is scrolled. A view is one scanned image: usually one side of a leaf.
\newcommand{\page}[1]{%
  \marginnote{\footnotesize\color{gray!70}\textsf{v.\,#1}}%
  \label{page:#1}\ignorespaces}

% The BnF's pencilled foliation, where the leaf carries it and a pass can read
% it: « 198r ». This is what the literature cites — « ff. 198r–208v » — and
% the one line that turns a citation into a view. Recorded, never inferred:
% a folio a pass cannot read is not written.
\newcommand{\folio}[1]{%
  \marginnote{\footnotesize\color{gray!50}\textsf{f.\,#1}}\ignorespaces}

% The modernised reading's coarser counterpart to \page{}: which views a
% section is drawn from.
\newcommand{\pagerange}[2]{%
  \marginnote{\footnotesize\color{gray!70}\textsf{v.\,#1--#2}}%
  \ignorespaces}

% Illegible: never guessed.
\newcommand{\ill}{\textcolor{red!60!black}{[\,\ldots\,]}}

% A doubtful reading: the word is given, and flagged as uncertain.
\newcommand{\uncertain}[1]{\uline{#1}}

% An editorial addition — a missing letter, an implied word.
\newcommand{\add}[1]{\textcolor{blue!50!black}{[#1]}}

% Struck out by the writer. What was crossed out is part of the document.
\newcommand{\struck}[1]{\sout{#1}}

% The transcriber's note, on the state of the page or the mathematics.
\newcommand{\note}[1]{\par\smallskip{\small\color{gray!60!black}\textit{[#1]}}\par\smallskip}

% A marginal note of hers, distinct from the body of the text.
\newcommand{\marginal}[1]{\par\smallskip\noindent\hspace{1em}%
  {\small\color{gray!60!black}#1}\par\smallskip}

% --- Title ------------------------------------------------------------------
% The title block is French, like the documents themselves.

\AddToHook{shipout/background}{%
  \ifx\grothWatermark\empty\else
    \begin{tikzpicture}[remember picture, overlay]
      \node[rotate=52, scale=3.4, align=center, text=gray!18,
            font=\sffamily\bfseries]
        at (current page.center) {\grothWatermark};
    \end{tikzpicture}%
  \fi}

\AtBeginDocument{%
  \begin{center}
    {\large\bfseries Manuscrits de mathématiciens (Gallica) ---
      \ifx\grothShelfmark\empty\grothFolder\else\grothShelfmark\fi}\\[2pt]
    {\itshape\grothFoldertitle}\\[4pt]
    {\small vues \grothFirst--\grothLast
      \ifx\grothBatch\empty\else\ (lot \grothBatch)\fi}\\[2pt]
    \ifx\grothDating\empty\else
      {\small\color{gray!60!black}Datation du catalogue : \grothDating}\\[2pt]
    \fi
    \ifx\grothWatermark\empty\else
      {\small\bfseries\color{red!45!gray}\grothWatermark}\\[2pt]
    \fi
    {\footnotesize\color{gray!60!black}Transcription établie d'après les images IIIF de Gallica.
     Source gallica.bnf.fr / Bibliothèque nationale de France.\\
     Les lectures incertaines sont soulignées, les illisibles notées [\,\ldots\,],
     les ajouts entre crochets ; \textsf{v.} numérote les vues de Gallica,
     \textsf{f.} la foliotation au crayon de la BnF.}
  \end{center}
  \vspace{1em}}

\endinput


\folder{latin-17859}
\shelfmark{Latin 17859}
\pages{1}{49}
\dating{1601-1700}
\watermark{Lecture automatique\\interprétation non vérifiée}
\foldertitle{Opuscules de François Maurolyco sur les mathématiques et l'astronomie — lecture modernisée du volume entier}

\begin{document}

\section*{Résumé}

\begin{resume}
Ces feuillets sont la copie, faite au XVII\textsuperscript{e} siècle d'après un
livre imprimé à Messine en 1558, de ce que Francesco Maurolico avait réuni sur la
sphère : la géométrie des cercles tracés sur une sphère, et son usage pour
calculer la position des astres. Maurolico dit lui-même, dans la lettre au
vice-roi de Sicile qui ouvre le volume, pourquoi il l'a fait : la science des
triangles sphériques, aussi nécessaire à la géométrie et à l'astronomie qu'elle
est négligée dans les écoles, ne pouvait rester dans cet abandon ; il a donc
restauré les \emph{Sphériques} des Grecs et y a joint ses propres travaux. La
copie n'est pas complète — les \emph{Sphériques} de Théodose et de Ménélas n'y
ont que leurs titres — et la reliure a déplacé quelques feuillets, que cette
lecture remet dans l'ordre du texte en le disant.

Le problème est celui de toute l'astronomie de position. Les astres sont vus sur
une sphère ; ce qu'on mesure (une hauteur, un lieu sur l'écliptique) et ce qu'on
veut connaître (une déclinaison, une ascension droite, l'heure, la durée du jour,
le point de l'horizon où le Soleil se lève) sont des côtés et des angles de
triangles dont les côtés sont des arcs de grands cercles. Il faut donc des
relations entre ces côtés et ces angles, et des tables pour calculer avec elles.

Le volume en donne les deux étages. D'abord la cosmographie et les textes grecs :
le \emph{Discours sur la sphère}, qui expose le monde géocentrique de 1558 ; la
\emph{Sphère en mouvement} d'Autolycus, les \emph{Habitations} de Théodose et les
\emph{Phénomènes} d'Euclide, où l'on apprend ce qui se lève en même temps, combien
de temps un signe du zodiaque met à se lever, pourquoi il y a sous le cercle
polaire un jour de vingt-quatre heures. Ensuite les outils : une table de sinus,
une table de tangentes (la « table féconde », attribuée à Regiomontanus) et une
table de sécantes (la « table bienfaisante », que Maurolico dit avoir faite à
son imitation), toutes pour un rayon de $100\,000$ ; huit problèmes entièrement
chiffrés — la déclinaison et l'ascension droite d'Aldebaran, l'amplitude du
Soleil levant à Messine, l'angle horaire du Soleil d'après sa hauteur ; et
surtout la fin du premier livre et le second livre entier des \emph{Sphériques}
de Maurolico, une soixantaine de propositions qui démontrent les relations dont le calcul a
besoin.

L'idée qui porte ces livres est de ramener chaque question à un triangle
rectangle, et chaque relation à une proportion entre sinus. On y trouve, sous des
formes et des noms qui ne sont pas les nôtres, le théorème de Ménélas, la règle des
sinus, la loi des cosinus sphérique écrite en sinus verses, les relations du
triangle rectangle que résument aujourd'hui les règles de Napier
($\cos c = \cos a\cos b$, $\cos A = \cos a\sin B$, $\sin b = \tan a\cot A$), le
théorème de Ptolémée, et un « très noble théorème », attribué à Ménélas, que
Maurolico démontre de trois manières : dans un triangle rectangle, le sinus de la
somme de l'hypoténuse et d'un côté est au sinus de leur différence comme
$1 + \cos A$ est à $1 - \cos A$, soit $\cot^2(A/2)$. Plusieurs propositions en tirent
une question d'extremum : l'angle $A$ étant fixé, la différence entre
l'hypoténuse et le côté est la plus grande quand leur somme vaut un quadrant, et
elle vaut alors $\arcsin\tan^2(A/2)$, valeur que le feuillet ne calcule pas. Un texte daté de 1530 explique enfin
pourquoi les règles des tables de Regiomontanus sont justes et comment la table
des sécantes permet de remplacer ses divisions par des multiplications.

Les énoncés sont presque tous justes, et cette lecture les a vérifiés, tables et
exemples chiffrés compris. Ce qui ne l'est pas tient à la copie plus qu'aux mathématiques :
des lettres de figure échangées, une colonne de soixantièmes qui a glissé de cinq
lignes, une tangente de $66^\circ 30'$ mal prise dans la table, une déclinaison de
$15^\circ 20'$ écrite pour $15^\circ 10'$, un cas oublié dans une réciproque ; on dit
chaque fois ce qui est vrai et, en note, ce que porte le feuillet. Le
\emph{Discours} et la note sur la visibilité de la Lune, eux, contiennent des
arguments que la physique a réfutés (l'immobilité de la Terre) ou qui ne tiennent
pas (un croissant visible à trois degrés du Soleil), et la lecture le dit.

Où cela mène : à la trigonométrie sphérique telle qu'on l'enseigne encore en
astronomie et en navigation — triangle sphérique, règle des sinus, loi des
cosinus, règles de Napier —, aux fonctions tangente et sécante et à l'histoire des tables
trigonométriques, et, du côté de l'astronomie, aux coordonnées équatoriales et
écliptiques, à la différence ascensionnelle, à l'amplitude et à l'angle horaire,
aux temps d'ascension des signes et aux levers héliaques.

\keywords{spherical trigonometry, right spherical triangle, Menelaus's theorem,
spherical law of sines, spherical law of cosines, versine, Napier's rules,
half-angle formula, sine addition formula, Ptolemy's theorem, trigonometric
tables, tangent table, secant table, declination, right ascension, ascensional
difference, rising amplitude, hour angle, oblique ascension, rising times of the
zodiacal signs, heliacal rising, spherical astronomy, geocentric cosmology,
trepidation}
\end{resume}

\section*{Ce que le volume contient}

\pagerange{1}{49}
Quarante-neuf vues d'un microfilm en niveaux de gris, chacune montrant le
volume ouvert : une page de gauche et une page de droite, que la
transcription a coupées et lues comme deux pages. La vue 1 (le plat, l'étiquette
de cote, et la garde où la Bibliothèque a certifié « Volume de 47 Feuillets.
7 Juin 1870 ») n'est pas transcrite ; la page de droite de la vue 49 est
blanche. Tout le reste est d'une seule main, une italique régulière, avec les
figures dans les marges.\footnote{Chaque page de droite des vues 2 à 48 porte en
haut, à l'encre et d'un module plus gros que le texte, un numéro : 1 à 47, un
par feuillet, sans lacune dans l'ordre présent de la reliure. La série et le
certificat de 1870 s'accordent sur quarante-sept feuillets. Les trois
transcriptions ne tiennent pourtant pas pour établi que ce soit la foliotation
de la Bibliothèque,
parce que c'est de l'encre et non le crayon fin d'un bibliothécaire, et que le
trait pourrait être celui du copiste ; aucune n'écrit de \emph{folio}, et cette
lecture ne cite donc les feuillets que par les vues de Gallica et, au besoin,
par ces « numéros à l'encre ». Le cachet rond « Bibliothèque impériale
MSS » est apposé dans les marges des vues 2, 20 et 49.}

C'est une copie, et elle le dit. La page de titre énumère le contenu d'un
recueil imprimé — les \emph{Sphériques} de Théodose « selon la tradition de
Maurolico de Messine, mathématicien », les \emph{Sphériques} de Ménélas, la
\emph{Sphère en mouvement} d'Autolycus, les \emph{Habitations} de Théodose, les
\emph{Phénomènes} d'Euclide « démontrés très brièvement », la « démonstration
et la pratique de trois tables, du sinus droit, féconde et bienfaisante,
appartenant aux triangles sphériques », un \emph{Compendium Mathematicæ} et le
\emph{Discours sur la sphère} de Maurolico — puis sa source : « Exscripta ex
typis excusso codice Messanæ. In freto Siculo a Petro Spira mense Augusto Anno
M.D.LVIII. »\footnote{« Copié sur l'exemplaire imprimé à Messine, sur le
détroit de Sicile, par Petrus Spira, au mois d'août de l'an 1558. » Sous cette
ligne, d'un module plus petit : « Bibliotheca I.U. Thuanorum. mense Iunio. Anno
M.DC.XXVIII. favore C.V.D.I.B. Haultini mense decembri mutuo concesso Anno
M.DC.XXVII. », une note de provenance des Thou (juin 1628), l'exemplaire ayant
été prêté en décembre 1627 par la faveur d'un « Haultinus » ; les lectures
« I.U. » et « Haultini » sont incertaines, le milieu du second mot étant
surchargé. En marge, d'une autre main, l'ancienne cote « Nouv. acq. lat.
1099 » ; au crayon, d'une main moderne, « Ce MS paraît incomplet », le
dernier mot mal lu (vue 2, page de droite ; la page de gauche est blanche).}
La datation du catalogue, 1601--1700, est reproduite en tête telle quelle. Le
volume donne lui-même un repère : sa toute dernière ligne (vue 49) omet le
\emph{Compendium} « parce que le P. Mersenne l'a publié », ce qui place la
copie après cette publication ; la date de celle-ci n'est pas sur les feuillets,
et cette lecture n'en infère aucune.

Le recueil copié n'est pas complet. Les \emph{Sphériques} de Théodose et de
Ménélas n'y ont que leurs titres, sur une page par ailleurs blanche, avec pour
Ménélas la mention « Alibi descripta » — copié ailleurs.\footnote{Vue 10, page de
gauche : le titre des trois livres de Théodose, « DIFFINITIONES. », la ligne
« Sed hi collati sunt cum veteri editione Veneta anni 1518. » (« ils ont été
collationnés avec l'ancienne édition vénitienne de 1518 »), puis « MENELAI
SPHÆRICA EX EIUSDEM TRADITIONE » et « Alibi descripta. » Ni définitions ni
texte ne suivent.} Ce qui est copié se range en neuf pièces :

\begin{itemize}
  \item deux lettres dédicatoires, l'une au vice-roi de Sicile (1558), l'autre
    à Charles Quint (1556), et l'\emph{Index} des travaux de Maurolico
    (vues 2--4) ;
  \item le \emph{Discours sur la sphère}, une cosmographie en prose
    (vues 4--10) ;
  \item trois textes grecs dans la « tradition » de Maurolico : Autolycus,
    \emph{De la sphère en mouvement} (vues 48 et 10--11), Théodose,
    \emph{Des habitations} (vues 11--12), Euclide, \emph{Phénomènes}
    (vues 13--15) ;
  \item deux notes sur les habitants de la Terre et sur la visibilité des
    astres près du Soleil (vues 15--16) ;
  \item trois tables — des sinus, des tangentes, des sécantes — et une table
    à trois arcs pour les déclinaisons et les ascensions (vues 16--20) ;
  \item huit problèmes résolus avec ces tables (vues 17--19 et 48--49) ;
  \item la fin d'un premier livre de trigonométrie sphérique, depuis la
    proposition XXIII, et un second livre entier, trente-deux propositions
    (vues 20--47) ;
  \item un commentaire des règles de Regiomontanus et la démonstration de la
    table des sécantes, daté d'août 1530 (vues 42--45 et 47) ;
  \item une dédicace à Octavio Spinola datée du 1\textsuperscript{er} juin 1557,
    et la ligne finale sur le \emph{Compendium} (vues 47 et 49).
\end{itemize}

\subsection*{L'ordre des feuillets, et celui de cette lecture}

La reliure a mis trois groupes de feuillets hors de l'ordre du texte, et les
transcriptions l'ont chaque fois montré par la continuité des phrases, en
gardant elles-mêmes l'ordre des vues.

\begin{itemize}
  \item La page de gauche de la vue 48 ouvre le livre d'Autolycus
    (définition, propositions 1 à 4) et s'arrête sur « et per primam huius
    puncta singula » ; la page de droite de la vue 10 commence par « suos
    seorsum parallelos, semper in alterutro hæmisphæriorum describent » et
    continue avec les propositions 5 à 12. La phrase se lit d'une traite.
  \item Les problèmes 1 à 6 se lisent sur les pages numérotées 17 (vue 18,
    droite), puis vue 19, gauche, puis 16 (vue 17, droite), puis vue 18,
    gauche : les deux feuillets numérotés 16 et 17 sont reliés à l'envers.
  \item La fin du second livre des \emph{Sphériques} court de la vue 41,
    gauche, à la vue 42, gauche, s'interrompt au milieu de la proposition XXVIII, et reprend
    à la vue 45, droite, qui répète les deux dernières lignes interrompues ; elle
    s'achève à la vue 47, gauche, sur une phrase en suspens, que continue la
    vue 42, droite. Le commentaire sur Regiomontanus qui occupe les vues 42,
    droite, à 45, gauche, s'achève à son tour en tête de la vue 47, droite.
    Les trois feuillets numérotés 41 à 43 appartiennent donc entre ceux
    numérotés 45 et 46.
\end{itemize}

Cette lecture suit l'ordre du texte là où ces notes l'établissent, et le dit
en note à chaque endroit. Elle fait en outre un rapprochement qui est le sien :
les problèmes « 7 » et « 8 » de la fin du volume (vues 48--49) continuent la
numérotation et reprennent les données (latitude de 38°, amplitude de
$30^\circ 24'$) du problème 6 de la vue 18, et sont lus à sa suite.\footnote{Aucune
note des transcriptions ne fait ce rapprochement. La transcription de la vue 18
dit la page blanche sous le problème 6 ; la vue 48, droite, s'ouvre sur « qui
quidem arcus est horizontis inter æquinoctialem, et solem in tali horizonte
orientem » (« lequel arc est celui de l'horizon entre l'équateur et le Soleil
se levant sur cet horizon »), ce qui définit l'amplitude du problème 6 et
pourrait achever sa dernière phrase ; ce n'est pas établi.} Pour le reste,
l'ordre de la reliure est gardé. Le texte laisse pourtant voir un autre ordre,
celui du livre imprimé : le commentaire de 1530 se termine par « Nunc
adnectemus his quædam veterum opuscula ad eandem sphæricam theoriam
spectantia », et la lettre à Spinola annonce qu'« une fois achevés les
\emph{Sphériques} de Théodose, de Ménélas et de ton Maurolico », viendront
« ce qui touche la sphère mobile et le premier mouvement ».\footnote{« Nous
joindrons maintenant à ceci quelques opuscules des anciens qui regardent la
même théorie sphérique » (vue 47, droite) ; « Absolutis itaque Theodosij
Menelai ac Maurolyci tui Sphæricis, subiungemus ea quæ ad sphæram mobilem
motumque primum pertinent » (même page). Dans le livre dont la copie est
tirée, les \emph{Sphériques} précédaient donc Autolycus, Théodose et Euclide ;
la copie, telle qu'elle est reliée, met le \emph{Discours} en tête.}

\section*{Les dédicaces et l'index des travaux}

\pagerange{2}{4}
La première lettre est adressée à Juan de la Cerda, duc de Medinaceli et
vice-roi de Sicile, et datée de Messine, août 1558.\footnote{« D. Iohanni Cerda
Methymnensium Duci ac Siciliæ Proregi Illustrissimo Maurolycus S.D. » ;
« Methymnensium Duci » est rendu par le titre castillan, ce qui est une
traduction et non une lecture. Une tache d'encre entre « M.D. » et « LVIII »
(vue 3, page de gauche).} Maurolico y raconte que ces petits livres des
\emph{Sphériques}, « corrigés au prix de beaucoup de veilles », étaient destinés
à Charles Quint trois ans plus tôt, et que le malheur des temps a empêché
l'envoi ; il prie le vice-roi — qu'il a vu récemment en armes sur le rivage
de Messine repousser la flotte turque qui ravageait la Calabre — de le
présenter à l'empereur. Il dit en une phrase ce qu'il a voulu faire, et c'est
la meilleure introduction au volume : la science des triangles sphériques,
aussi nécessaire à la géométrie et à l'astronomie qu'elle est négligée dans les
écoles, ne pouvait rester dans cet abandon ; il a donc restauré, « à partir
d'exemplaires très corrompus », les \emph{Sphériques} de Théodose et de
Ménélas, et y a joint « quelques veilles » de lui et d'autres pièces sur le
même sujet.\footnote{« Sphæralium quidem triangulorum scientia, quantum erat
Geometricis et Astronomicis Studiis necessaria tantum in Gymnasijs hactenus
neglecta iacuit. Nec erat ulterius tantus Academiæ speculationum tolerandus
defectus. Itaque quoad potui conatus sum ut Theodosii et Menelai sphærica ex
corruptissimis exemplaribus instaurata in lucem prodirent. His nonnullas
lucubrationes meas et alia quædam ad idem negotium spectantia subiunxi. » —
« La science des triangles sphériques, autant elle était nécessaire aux études
de géométrie et d'astronomie, autant elle est restée jusqu'ici négligée dans
les écoles ; et il ne fallait pas tolérer plus longtemps un tel manque dans les
spéculations de l'Académie. J'ai donc tâché, autant que je l'ai pu, que les
\emph{Sphériques} de Théodose et de Ménélas, restaurés d'après des exemplaires
très corrompus, paraissent au jour. J'y ai joint quelques-uns de mes travaux et
d'autres choses touchant la même matière. »}

Suit l'« Index lucubrationum Maurolyci », en deux parties. Les œuvres
d'autrui qu'il a éditées ou corrigées (« Aliena ») : les \emph{Éléments}
d'Euclide, « les erreurs des interprètes discutées », avec des propositions
ajoutées sur les solides réguliers ; les \emph{Sphériques} de Théodose et de
Ménélas ; quatre livres des \emph{Coniques} d'Apollonius ; Serenus ; les
traités d'Archimède sur la mesure du cercle, la sphère et le cylindre, les
isopérimètres, l'équilibre des plans, la quadrature de la parabole, les
sphéroïdes et conoïdes, les spirales ; l'arithmétique de Jordanus ; les
perspectives de Roger Bacon et de « Iohannes » ; les miroirs ardents de
Ptolémée ; Autolycus, Théodose et les \emph{Phénomènes} ; les \emph{Problèmes
mécaniques} d'Aristote.\footnote{Vue 3, pages de gauche et de droite (numéro à
l'encre 2). « errori[bu]s » : une tache couvre la fin du mot ; « adnotatibus »
porte un signe d'abréviation, pour « adnotationibus » ; « facilitas », à la
fin de l'article sur les spirales, est une lecture incertaine ;
« fabricatioq; » abrège « fabricationeque ».} Ses propres travaux (« Nostra
vero sunt ») : des prologues sur la division des arts ; une \emph{Arithmétique
spéculative} en deux livres, sur les formes des nombres et sur les grandeurs
rationnelles et irrationnelles ; des démonstrations des règles de fausse
position et de la « chose » réduites à quatre préceptes ; « deux petits livres
de \emph{Sphériques}, où beaucoup de choses négligées par Ménélas sont
suppléées » — ce sont, selon toute apparence, les deux livres dont ce volume
garde la fin du premier et le second entier ; une sphère mobile ; une
cosmographie dédiée jadis à Pietro Bembo ; les livres V et VI des
\emph{Coniques} ; un traité des solides réguliers qui remplissent l'espace, « où
l'ignorance d'Averroès apparaît manifestement » ; quatre livres sur l'égalité
des moments, dont le dernier sur les centres de gravité des solides omis par
Archimède ; les instruments ; trois livres sur les lignes horaires ; les
\emph{Photismi} sur la lumière et l'ombre, et les \emph{Diaphana} en trois livres
(corps transparents, arc-en-ciel, structure de l'œil et lunettes) ; des
questions arithmétiques, géométriques et astronomiques ; des annotations et des
canons de tables ; enfin le \emph{Compendium Mathematicæ}.\footnote{L'article
sur les lignes horaires contient un énoncé : les lignes des heures comptées
depuis midi coupent une certaine courbe aux points où celle-ci est touchée par
les lignes des heures comptées depuis le coucher et le lever, et cette courbe
est un cercle ou une section conique. Il est donné sans démonstration, et
cette lecture ne l'examine pas. Dans l'article des annotations, un mot d'environ
cinq lettres, à peu près « sioda », n'est pas lu. L'identification des deux
livres de \emph{Sphériques} de l'index avec ceux du volume est de cette
lecture.}

La seconde lettre, à Charles Quint, est datée de Messine, juillet 1556 — ou
d'une année voisine, le dernier chiffre étant surchargé.\footnote{« AD
C[arol]um V. Imperatorem maximum Francisci Maurolyci Siculi Abbatis Epistola » ;
le milieu du nom est effacé le long d'un pli ; la date se lit « M.D.LVI »
suivi d'un ou deux jambages surchargés, non lus. Un mot biffé, « tantopere »
(vue 4, page de gauche).} Elle justifie l'offrande d'un livre de sphères à un
souverain par la sphère elle-même : la figure la plus simple, la plus égale, la
plus forte et la plus capable, image de l'archétype, convient à qui gouverne la
sphère terrestre ; elle nomme ceux qui ont encouragé l'entreprise, Juan de Vega,
ancien vice-roi, et Simeone Ventimiglia, marquis de Geraci, alors stratège de
Messine. Un poème de douze vers, « Eidem Augusto Dicatum carmen », la
suit.\footnote{Deux colonnes de distiques. Au premier vers, « nuclatum » et
« viverat » sont des lectures incertaines ; les fins des six vers de la colonne
de droite passent dans le pli de la reliure et ne se voient pas. Le poème n'est
donc pas traduit ici ; son dernier vers lisible partage l'empire entre Jupiter,
qui a le ciel, et César, qui a « notre sphère ».}

\section*{Le Discours sur la sphère}

\pagerange{4}{10}
« Maurolyci Messanensis de Sphæra Sermo » est une cosmographie complète en
quelques pages, sans figures ni démonstrations : ce que doit savoir, en 1558,
celui qui va lire des traités de la sphère. On le suit ici dans son ordre, en
disant ce qui reste vrai et ce qui ne l'est plus.\footnote{Le \emph{Discours}
commence à la page de droite de la vue 4 (numéro à l'encre 3) et s'achève sur
« FINIS. » à la page de gauche de la vue 10. Les pages s'enchaînent sans lacune ;
chaque transcription de page note les mots qui font la jonction.}

\subsection*{La sphère, figure parfaite}

Le cercle parmi les figures planes, la sphère parmi les solides, l'emportent
par six propriétés : la simplicité (une seule ligne, une seule surface les
bornent), la similitude des parties (tous leurs arcs, toutes leurs calottes
sont également courbés), l'égalité (le centre est à égale distance du bord, et
tous les diamètres sont égaux), l'identité du lieu (tournant autour du centre,
la figure occupe toujours le même lieu), la solidité, et la capacité (à
périmètre égal le cercle, à surface égale la sphère enferment le plus). Il
ajoute que chacune suffirait à définir la figure.\footnote{« Quarum sex
proprietatum quælibet satis esset ad figuræ diffinitionem : cum singula solis
conveniant. » Au mot « æqualite », ainsi écrit pour « æqualitate », il n'y a
pas de signe d'abréviation visible (vue 4, droite).}

C'est vrai de quatre d'entre elles et faux de deux. Une courbe fermée du plan
dont la courbure est constante est un cercle ; une courbe invariante par toutes
les rotations autour d'un point est un cercle ; l'équidistance au centre est la
définition même ; et parmi les courbes fermées de longueur donnée, le cercle
seul enferme l'aire maximale — c'est l'inégalité isopérimétrique, dont les
démonstrations complètes sont bien postérieures au feuillet.\footnote{Le
feuillet ne démontre rien de cela et ne prétend pas le faire ; il renvoie à
Euclide, Ptolémée, Théodose, Ménélas et Archimède pour les propriétés de la
sphère. L'égalité de tous les diamètres, qu'il range sous l'« égalité », ne
caractérise pas le cercle : les courbes de largeur constante, dont le triangle
de Reuleaux, ont tous leurs « diamètres » égaux sans être des cercles ; ce
n'est vrai que si l'on entend par diamètre une corde passant par un centre
fixe.} En revanche, être bornée par une seule ligne ne caractérise rien — toute
courbe fermée simple l'est —, et la « solidité » n'est pas une propriété
géométrique. Une septième propriété, qu'on ne puisse assigner ni commencement
ni fin, fait passer à la cosmologie : Dieu a donné cette forme au monde pour
qu'il imite un archétype sans commencement ni fin.\footnote{« Quo fit ut non
temere neque Sed consulto […] credamus factum ut mundus huiusmodi formam
contraxerit » : « Sed » est écrit dans l'interligne au-dessus de « neque », qui
n'est pas biffé. La phrase finale, « la sphère est engendrée par un demi-cercle
tournant autour de son diamètre fixe », porte « semicirculo » écrit au-dessus
de « diametro » biffé (vues 4 droite et 5 gauche).}

\subsection*{La Terre ronde, au centre, immobile}

Le ciel est sphérique et tourne circulairement : une étoile garde toujours la
même grandeur apparente, et les phases de la Lune ne s'expliqueraient pas sans
un astre sphérique. La Terre et la mer forment un globe : les astres se lèvent
plus tôt pour qui est plus à l'est ; en marchant sur un méridien, la hauteur
méridienne d'un astre et la longueur de son arc diurne varient ; et comme ces
écarts sont proportionnels aux distances parcourues, la surface est ronde. On
voit du haut des collines plus de mer que du rivage ; l'ombre de la Terre sur la
Lune éclipsée est ronde ; l'équilibre des graves, qui tendent tous vers le
centre, exige une surface dont tous les points soient à égale distance du
centre ; les montagnes ne font pas plus à la rondeur du globe que les aspérités
d'une grosse boule de pierre.\footnote{Dans la phrase sur la proportionnalité,
un mot de six lettres environ, à peu près « tetium », n'est pas lu (vue 5,
gauche) : « Cumque [\ldots] ortivum intervalla et crementa sint interiectis terræ
spatijs proportionalia, circularem rotunditatem esse necesse est. » Vue 5,
droite : « iam » est offert pour un mot qui se lit « ium » ; « quieta » est
surchargé ; « stultitia » se lit plutôt « Shelhitia » au trait et n'est offert
que par le parallèle avec « Cedat antiquorum ignorantia ».} Il écarte ceux qui
donnaient à la terre et à la mer deux sphères ou deux centres distincts : le
globe entier, terre et eaux ensemble, a pour centre de gravité le centre du
monde.

La Terre est au centre, puisqu'une étoile paraît de même grandeur de partout,
qu'on voit toujours la moitié du zodiaque, que les extrémités des ombres
équinoxiales sont en ligne droite, qu'il n'y a d'éclipse de Lune qu'à
l'opposition, et qu'un astre se couche quand son opposé se lève ; et elle n'est
qu'un point au regard du firmament, de sorte que la distance de l'œil au centre
de la Terre ne change pas sensiblement les observations.

Elle est immobile. Un mouvement droit lui ferait quitter le centre ; une
rotation autour d'un axe autre que celui du monde ferait varier la hauteur du
pôle en un même lieu, ce qu'on n'observe pas ; une rotation autour de l'axe du
monde laisserait en arrière, vers l'occident, les nuages et les oiseaux, et une
pierre lancée en l'air ne retomberait pas au même endroit. La deuxième raison
est juste : une rotation diurne autour d'un autre axe se verrait. La troisième
ne l'est pas. Les nuages, les oiseaux et la pierre partagent le mouvement de la
Terre et le conservent ; une rotation uniforme du globe autour de l'axe du monde
ne produit, à l'échelle de ces expériences, aucun des effets invoqués, et c'est
bien la Terre qui tourne.\footnote{Le principe d'inertie et la relativité du
mouvement uniforme, qui réfutent cet argument, sont formulés au siècle suivant,
après le feuillet ; les effets réels de la rotation à cette échelle (ceux de la
force de Coriolis) se mesurent en centimètres au plus, sans commune mesure avec
ceux que l'argument annonce. Le feuillet dit en outre que le mouvement « est
contraire à un corps grave et en repos », ce qui est une thèse de physique
aristotélicienne, non un fait.}

\subsection*{Les deux mouvements et l'ordre des sphères}

Le ciel a deux mouvements : le premier, d'orient en occident, qui fait lever et
coucher tous les astres ; le second, en sens contraire, visible dans le
déplacement du Soleil et l'inégale vitesse de la Lune et des planètes. Chaque
mouvement demandant son orbe, il faut au moins huit cieux ; un neuvième pour le
lent mouvement des fixes vers l'orient que Ptolémée a perçu, et que Thābit a
changé en « trépidation » ; un dixième depuis qu'Alphonse a ajouté à la
trépidation un mouvement en longitude. L'ordre des sphères se tire de la lenteur
du mouvement propre (la huitième sphère la plus haute, la Lune la plus basse),
des éclipses (la Lune cache le Soleil), de la parallaxe, et de la convenance :
le Soleil au milieu des planètes, Mars et Vénus à ses côtés.\footnote{« Dispositio
autem et ordo sphærarum sumitur a quantite secundarij motus » : « quantite »
porte un signe d'abréviation, pour « quantitate » (vue 6, gauche).} Tout cela
est la cosmologie géocentrique reçue en 1558, que le feuillet expose sans
discussion et que l'astronomie héliocentrique a remplacée.

\subsection*{Les cercles de la sphère, les horizons, les ascensions}

Suit le vocabulaire qu'utilisent tous les textes du volume, et qui est encore
le nôtre : l'équateur, grand cercle perpendiculaire à l'axe du mouvement diurne ;
l'écliptique (le « zodiaque »), qui le coupe obliquement et touche les deux
tropiques ; les colures, par les pôles du monde et, l'un, par les pôles de
l'écliptique et les solstices, l'autre par les équinoxes ; l'horizon, dont le
pôle est le zénith, droit s'il passe par les pôles du monde, oblique sinon ; le
méridien, le premier vertical ; les cercles de latitude (par les pôles de
l'écliptique), de déclinaison (par les pôles de l'équateur) et de hauteur (par
le zénith). La latitude d'un lieu est l'arc du méridien entre le zénith et
l'équateur, égal à la hauteur du pôle ; sa longitude se compte sur l'équateur
depuis le méridien de l'extrême occident.\footnote{« Circulus altitudinis per
polos \emph{altitudinis} horizontis » : le premier « altitudinis » est biffé
(vue 6, gauche).}

L'horizon droit coupe tous les parallèles en deux moitiés : tout astre s'y lève
et s'y couche, et l'équinoxe y est perpétuel. L'horizon oblique touche deux
parallèles égaux, entre lesquels et les pôles les astres ne se couchent jamais
ou ne se lèvent jamais ; il coupe les autres inégalement, les arcs diurnes étant
plus grands du côté du pôle visible, et deux parallèles symétriques par rapport
à l'équateur ont l'arc diurne de l'un égal à l'arc nocturne de l'autre. Entre les
tropiques, le Soleil passe deux fois l'an au zénith et l'écliptique est deux
fois par jour perpendiculaire à l'horizon. Là où le pôle de l'horizon est sur le
cercle arctique, le plus long jour vaut vingt-quatre heures, avec « un instant »
pour nuit ; plus près du pôle, le Soleil reste levé tant qu'il parcourt un arc
d'écliptique retranché, « autant de jours que de degrés, autant de mois que de
signes » ; au pôle, six mois de jour et six de nuit.\footnote{« instans » est
offert pour un mot qui se lit à peu près « insterps », par le parallèle « et
instans pro nocte » de la vue 12 ; « circulus » est écrit pour « circulum ».
« Autant de jours que de degrés » suppose le Soleil avançant d'un degré par
jour, ce qui est vrai à un peu plus d'un pour cent près (vue 6, droite).}

Les parties de l'écliptique ne montent pas uniformément, puisque l'écliptique
tourne autour d'un axe qui n'est pas le sien. L'\emph{ascension} d'un arc
d'écliptique est l'arc d'équateur qui se lève avec lui ; elle est droite sur
l'horizon droit, oblique sur l'horizon oblique ; la \emph{descension} est
l'arc d'équateur qui se couche avec lui. Les énoncés du feuillet sont exacts :
sur l'horizon droit, quatre arcs également distants des équinoxes ont même
ascension, et un arc plus proche des solstices en a une plus grande ; deux arcs
diamétralement opposés se lèvent quand l'autre se couche ; pour tout arc,
ascension et descension obliques ont une somme indépendante du lieu (le double
de l'ascension droite, puisque la différence ascensionnelle s'ajoute à l'une et
se retranche de l'autre) ; un arc symétrique par rapport à un solstice a même
ascension et même descension ; dans nos latitudes nord, un arc de la moitié
ascendante de l'écliptique (du Capricorne aux Gémeaux) a l'ascension plus petite
et la descension plus grande. Sur l'horizon droit deux demi-cercles
quelconques de l'écliptique se lèvent en temps égaux ; sous le cercle arctique,
le demi-cercle ascendant se lève en un instant et se couche en vingt-quatre
heures.\footnote{La somme « conficiunt eumdem arcum » est lue ici « le même arc
en tout lieu » ; c'est l'interprétation qui rend l'énoncé vrai. Avec
$\alpha$ l'ascension droite et $\Delta$ la différence ascensionnelle d'un
point, l'ascension oblique est $\alpha-\Delta$ et la descension oblique
$\alpha+\Delta$ (latitude nord) : pour un arc, la somme vaut $2\,\delta\alpha$.}

\subsection*{Levers et couchers relatifs au Soleil}

Un astre qui se lève au lever du Soleil fait un lever matinal, s'il se couche à
ce moment un coucher matinal ; au coucher du Soleil, les deux sont vespéraux ;
pour une étoile de l'écliptique, le lever matinal et le coucher vespéral tombent
le même jour, de même le lever vespéral et le coucher matinal.\footnote{« stella
oriens fuit ortum » : le sens demande « facit » (vue 7, gauche).} Les astres plus
lents que le Soleil cessent d'être vus le soir quand il les rejoint (dernière
apparition vespérale), puis reparaissent le matin (première apparition
matinale) ; la Lune, plus rapide, fait l'inverse : dernière apparition le
matin, première le soir.\footnote{Le feuillet écrit « postrema fulsio
vespertina » là où le parallèle avec ce qui précède demande « prima » ;
« vespertina » est repassé à l'encre.} Il ajoute que la Lune a « parfois » été vue
les deux fois le même jour. Cela exigerait qu'elle soit assez loin du Soleil
pour être visible — au moins sept degrés environ — le matin puis le soir d'un
même jour, soit un déplacement relatif d'au moins quatorze degrés en moins
d'un jour, à la limite extrême de ce que permet son mouvement (douze à quinze
degrés par jour) : l'observation serait tout au plus exceptionnelle, et rien sur
le feuillet ne l'établit.\footnote{La limite d'environ 7° d'élongation au-dessous
de laquelle le croissant n'est pas visible (limite de Danjon) est une mesure du
XX\textsuperscript{e} siècle. Voir aussi, plus bas, la section sur la visibilité des astres, où le feuillet
tire trois degrés d'un calcul qui ne tient pas.}
Vénus et Mercure, tantôt plus rapides tantôt plus lents que le Soleil, ont les
deux sortes d'apparitions ; les étoiles proches du pôle visible, à cause de la
longueur de leur arc diurne, peuvent être vues le matin avant d'avoir cessé de
l'être le soir. Les signes « physiques » sont sous-tendus par les côtés de
l'hexagone inscrit ; divisés en deux, ils donnent les douze signes de trente
degrés, comme l'année a douze mois de trente jours.

\subsection*{Le jour, les heures, les habitants, les climats}

Le jour naturel est le temps d'une révolution de l'équateur augmentée de l'arc
correspondant au mouvement du Soleil pendant ce temps ; ces jours sont inégaux
(les jours « vulgaires » ou apparents) ; en répartissant également sur l'année
le cercle entier que font leurs suppléments, on obtient les jours égaux,
« astronomiques », qui servent au calcul des mouvements. Leur différence est
l'« équation des jours » — notre équation du temps.\footnote{« Cumque in toto anno
additamenta additamenta diurna collecta perficiant circulum » : « additamenta »
est écrit deux fois de suite ; « vulg. » en fin de ligne, puis « Inæquales »
(vue 7, droite).} Le feuillet la dit « à soustraire » du temps vulgaire ; c'est
vrai si on la compte à partir de sa valeur extrême, faute de quoi elle change de
signe au cours de l'année, entre environ $-14$ et $+16$ minutes.

L'heure équinoxiale est la vingt-quatrième partie du jour ; l'heure temporelle,
la douzième partie du jour ou de la nuit. Les périèques habitent le même
parallèle ; les antèques des parallèles égaux et opposés ; les antipodes, qui
sont en outre diamétralement opposés, ont le même horizon et voient chacun ce
que l'autre ne voit pas.\footnote{« dispõo » abrège « dispositio » ;
« stantes » est une lecture incertaine (vue 7, droite).} Le jour se compte à
partir du méridien et non de l'horizon, parce que deux méridiens se succèdent
toujours au même intervalle de temps, ce que ne font deux horizons que pour des
périèques. Ptolémée distingue les climats par des accroissements d'une
demi-heure du plus long jour, les parallèles par des quarts d'heure.

\subsection*{Les modèles des planètes}

L'inégalité des mouvements planétaires est « sauvée » par des excentriques et
des épicycles, qui rendent compte aussi des stations, des rétrogradations et des
latitudes. Le feuillet résume les modèles de Ptolémée : pour le Soleil, un
excentrique sans épicycle ou, ce qui revient au même, un concentrique portant
un épicycle ; pour les planètes supérieures et Vénus, un équant situé au-delà du
centre du déférent d'autant que celui-ci est au-delà du centre du monde ; pour
Mercure, un centre du déférent tournant sur un petit cercle ; pour la Lune, un
centre du déférent tournant en sens contraire de l'épicycle, et l'apogée moyen
de l'épicycle repéré depuis le point le plus bas de ce petit cercle. Les
planètes supérieures sont écartées de l'apogée de leur épicycle autant que le
Soleil l'est de l'épicycle, et les inférieures ont le même mouvement moyen que
le Soleil ; les équations se prennent dans les tables par le centre et
l'argument, corrigées par les minutes proportionnelles.\footnote{« per quos
salvantur salvantur » : le verbe est écrit deux fois (vue 8, gauche) ; « autem
autem » en fin et début de ligne (vue 8, droite) ; les deux dernières lignes de
la vue 7, droite, sont d'une encre plus pâle ; un passage biffé à la vue 8,
« remisso motu directo, redit ad directionem », est ensuite réécrit à sa place.}

Les éclipses ont lieu aux nœuds ou près d'eux, entre des limites où elles sont
possibles et d'autres, plus étroites, où elles sont nécessaires ; pour le Soleil,
tout le calcul se fait sur les lieux vus, corrigés de la parallaxe. Les latitudes
des planètes se composent de l'inclinaison du déférent et de celles de
l'épicycle ; pour Vénus et Mercure, de la « déviation » du déférent, de
l'« inclinaison » et de la « réflexion » de l'épicycle, avec les règles de
signe que le feuillet détaille.\footnote{Dans la description de l'épicycle des
planètes supérieures, un mot avant « dictum diametrum », à peu près « sontim »,
n'est pas lu (vue 8, droite). À la vue 9, gauche, une phrase entière sur
Vénus et Mercure est biffée puis réécrite. La dernière phrase de la section,
qui tire la supériorité de Vénus sur Mercure de ce que « le septentrional
précède en dignité l'austral », n'est pas un argument astronomique.} Tout cela
est l'astronomie des \emph{Tables alphonsines}, remplacée depuis par les
orbites elliptiques et la gravitation ; on ne le transpose pas en notations
modernes, parce qu'il n'y a pas d'énoncé moderne qui lui corresponde.

\subsection*{La huitième sphère}

La huitième sphère, outre le mouvement diurne, en a deux : la
« trépidation », par laquelle les extrémités du diamètre qui joint le début du
Bélier et celui de la Balance de la huitième sphère parcourent deux petits
cercles centrés sur les mêmes points de la neuvième, en 7000 ans ; et le
mouvement en longitude, dû à la neuvième sphère, en 49\,000 ans. L'équation de
trépidation monte au plus à 9°, le rayon du petit cercle ; c'est, dit-il, le
système d'Alphonse que suit l'académie commune des astronomes. Le mouvement réel
auquel correspondent ces constructions est la précession des équinoxes, qui est
uniforme à l'échelle historique, d'environ $50''$ par an, soit un tour en
quelque 26\,000 ans ; il n'y a pas de trépidation.

Suivent un éloge du Soleil, qui règle les planètes et donne au temps ses degrés,
ses mois et ses années ; les raisons des noms des constellations (la forme, la
Flèche ; l'influence, l'Autel qui fait les prêtres ; la nature, le Scorpion ;
l'histoire, Hercule et Persée ; la faveur des princes, la Chevelure de
Bérénice), et leurs variantes selon Hygin, Ptolémée et Alphonse ; et la division
des douze maisons, que Regiomontanus fait sur l'équateur, Campanus sur le
premier vertical, et qu'il vaudrait mieux, selon lui, rapporter à
l'écliptique.\footnote{« Quod ab illis Aquila ab \emph{illis} his vultur
volans » : « his » est écrit au-dessus de « illis » biffé (vue 9, droite).}

\subsection*{Les rencontres de deux mobiles}

Un seul passage du \emph{Discours} est un théorème, et il est juste. Deux mobiles
à vitesses constantes $v_1 > v_2$, séparés par une distance $d$, se rejoignent
après le temps $d/(v_1-v_2)$ s'ils vont dans le même sens, $d/(v_1+v_2)$ s'ils vont
en sens contraires. Sur un cercle, si les vitesses angulaires sont dans le
rapport $p:q$ en plus petits termes (le feuillet dit « les plus petits nombres
représentant la proportion des mouvements »), les rencontres ont lieu en
$p-q$ points également espacés dans le premier cas, $p+q$ dans le second, et
toujours aux mêmes points.\footnote{En effet, entre deux rencontres le premier
mobile parcourt $p/(p-q)$ de tour (resp. $p/(p+q)$), et comme $p$ est premier
avec $p\mp q$, les rencontres successives visitent exactement $p\mp q$ points
équidistants. L'hypothèse que $p$ et $q$ sont premiers entre eux est celle des
« plus petits termes ».} Si les vitesses sont incommensurables, les points de
rencontre ne se répètent jamais — ils sont même denses sur le cercle, ce que le
feuillet ne dit pas. D'où sa conclusion : si les mouvements célestes sont
incommensurables, les platoniciens se trompent en croyant que le ciel revient
un jour à la même disposition.

\subsection*{Comment chaque grandeur se trouve}

Le \emph{Discours} se termine par l'inventaire des observations qui fixent les
paramètres : l'obliquité par les hauteurs méridiennes du Soleil ; l'entrée du
Soleil à l'équinoxe par sa hauteur méridienne ; la longueur de l'année par deux
observations ; l'excentricité et l'apogée du Soleil par trois lieux ; le
mouvement de la Lune et de ses nœuds par des couples d'éclipses ; le diamètre de
l'épicycle lunaire par trois éclipses ; la distance de la Lune par la
parallaxe, et de là, par une éclipse de Soleil et les diamètres apparents, les
autres distances et les vraies grandeurs ; les révolutions des planètes par des
retours au même lieu et à la même distance du Soleil ; l'inclinaison du
déférent lunaire par la hauteur méridienne de la Lune quand ses nœuds sont aux
équinoxes. « Motus instauratis observationibus rectificantur, et radix
statuitur. FINIS. »\footnote{« Les mouvements sont rectifiés par des observations
renouvelées, et l'on fixe la racine » (la position à une époque donnée).
« radix » se lit à peu près « ruedix », sans certitude (vue 10, gauche).}

\section*{Autolycus, De la sphère en mouvement}

\pagerange{48}{48}
Le livre d'Autolycus commence à la vue 48 et se poursuit à la vue 10 ; on le
lit ici dans cet ordre.\footnote{La page de gauche de la vue 48 porte le titre
« Autolyci de sphæra quæ movetur ex traditione Maurolyci. Liber. », la
définition et les propositions 1 à 4, et s'arrête sur « et per primam huius
puncta singula » ; la page de droite de la vue 10 (numéro à l'encre 9) s'ouvre
sur « suos seorsum parallelos, semper in alterutro hæmisphæriorum describent »
et continue avec les propositions 5 à 12. Le raccord est établi par les
transcriptions des deux lots, qui gardent l'ordre des vues ; le titre de la
vue 2 place de même Autolycus après Ménélas, et Théodose note plus loin que sa
première proposition « est à peu près la quatrième de la sphère
d'Autolycus ».} C'est la cinématique de la rotation uniforme d'une sphère
autour d'un axe, rapportée à un grand cercle fixe, l'horizon. Maurolico donne
chaque énoncé suivi d'une preuve de deux ou trois lignes renvoyant aux
\emph{Sphériques} de Théodose.

\emph{Définition.} Des points d'une sphère se meuvent uniformément s'ils
parcourent en temps égaux des arcs égaux et semblables ; un point mû
uniformément sur une ligne parcourt des longueurs proportionnelles aux
temps.\footnote{« eandem habebit rationem tempus tempus » : « tempus » en fin de
ligne puis en début de ligne, sans « ad ». « Proposito 3. » est écrit ainsi
dans la marge ; la fin d'« hæmisphæriorum » est surchargée (vue 48, gauche).}

\begin{enumerate}
  \item[1.] Si la sphère tourne uniformément autour de son axe, tout point
    autre que les pôles décrit un cercle perpendiculaire à l'axe, ayant les
    mêmes pôles que la sphère.
  \item[2.] Tous les points parcourent en un même temps des arcs semblables de
    leurs parallèles (c'est-à-dire des arcs de même angle au centre).
  \item[3.] Réciproquement, les arcs de parallèles parcourus en un même temps
    sont semblables.
  \item[4.] Si l'horizon fixe est perpendiculaire à l'axe, aucun point ne se
    lève ni ne se couche : chacun décrit son parallèle dans l'un des deux
    hémisphères.
\end{enumerate}

\subsection*{Suite du livre}

\pagerange{10}{11}
\begin{enumerate}
  \item[5.] Si l'horizon passe par les pôles, tout point se lève et se couche,
    et reste autant de temps au-dessus qu'au-dessous : l'horizon coupe chaque
    parallèle en deux moitiés.\footnote{Dans les renvois, « per 20. primi
    Sphæricorum Theodosij », le 2 de cette main est un petit trait en forme de
    z, distinct du 1 (vue 10, droite).}
  \item[6.] Un horizon oblique touche deux parallèles égaux ; celui qui est du
    côté du pôle visible est toujours visible, l'autre toujours caché.
  \item[7.] Les parallèles se lèvent et se couchent toujours aux mêmes points
    de l'horizon, sous la même inclinaison.
  \item[8.] Tout grand cercle tangent aux deux parallèles que touche l'horizon
    coïncide avec l'horizon à un instant de la révolution.\footnote{« in sphæra
    in sphæra » est écrit deux fois.}
  \item[9.] De deux points qui se lèvent ensemble, le plus proche du pôle
    visible se couche le plus tard ; de deux points qui se couchent ensemble,
    le plus proche du pôle visible s'est levé le premier.
  \item[10.] Un grand cercle passant par les pôles est deux fois perpendiculaire
    à l'horizon oblique à chaque révolution.
  \item[11.] Un grand cercle tangent à des parallèles au moins égaux à ceux que
    touche l'horizon se lève et se couche sur tout l'arc d'horizon compris
    entre les parallèles qu'il touche.\footnote{« ortus et occasus
    \emph{fiet} facit » : « fiet » biffé (vue 11, gauche).}
  \item[12.] Si un cercle fixe coupe toujours en deux un cercle entraîné, et que
    ni l'un ni l'autre n'est perpendiculaire à l'axe ni ne passe par les pôles,
    ce sont deux grands cercles.
\end{enumerate}

Les énoncés sont exacts. La preuve du dernier procède par cas sans justifier le
premier : que deux petits cercles ne puissent se couper toujours en deux se
tire, en termes modernes, de ce qu'un plan ne coupe un cercle en deux moitiés
que s'il passe par son centre, ce que le mouvement du cercle entraîné rend
impossible à tout instant.\footnote{Le premier cas est écrit « Nam si uterque sit
circulus \emph{m} maior », après un début de mot biffé ; le raisonnement demande
« minor ». « ad ad rectos » est répété à la ligne (vue 11, gauche).}

\section*{Théodose, Des habitations}

\pagerange{11}{12}
« Theodosii De Habitationibus ex traditione Maurolyci » tire les conséquences
de la cinématique précédente pour des observateurs placés en divers lieux. Les
douze propositions, dans nos termes :

\begin{enumerate}
  \item[1.] Au pôle nord, le même hémisphère est toujours visible ; aucun astre
    ne se lève ni ne se couche (Autolycus 4).
  \item[2.] À l'équateur, tout astre se lève et se couche, et reste douze heures
    au-dessus de l'horizon (Autolycus 5).\footnote{Sur les vues 11, droite, et
    12, « Propositio » et son numéro sont écrits dans la marge de droite.}
  \item[3.] Entre les tropiques, l'écliptique devient perpendiculaire à
    l'horizon deux fois par jour ; sous un tropique, une seule fois, quand le
    point solsticial passe au zénith.
  \item[4.] Là où le zénith est éloigné du pôle autant que le tropique l'est de
    l'équateur — sur le cercle polaire —, six signes se lèvent à la fois, en un
    instant, une fois par révolution : l'écliptique se confond alors avec
    l'horizon.
  \item[5.] À l'équateur, quand les points solsticiaux sont sur l'horizon,
    l'écliptique passe par le zénith, et le méridien coupe en deux sa moitié
    visible.
  \item[6.] À l'équateur, tous les demi-cercles de l'écliptique se lèvent en des
    temps égaux.\footnote{L'énoncé s'arrête au bas de la vue 11, droite, sur
    « similiter etiam », et la page suivante, verso du même feuillet, commence
    avec la proposition 7 : la fin de l'énoncé et la démonstration ne sont pas
    écrites. On donne ce que l'énoncé écrit dit ; la suite, qui porterait
    vraisemblablement sur les couchers, n'est pas suppléée.}
  \item[7.] Sous un même parallèle, un astre se lève et se couche plus tôt au
    lieu le plus oriental, du même intervalle de temps.
  \item[8.] Sous un même méridien, un astre de déclinaison nord reste plus
    longtemps sur l'horizon pour l'observateur le plus au nord, et se lève plus
    tôt pour lui d'autant qu'il se couche plus tard ; pour un astre de
    déclinaison sud, c'est l'inverse.
  \item[9.] Pour deux lieux qui ne sont ni sur un même parallèle ni sur un même
    méridien, les arcs diurnes restent inégaux comme en 8, et les levers et
    couchers ne sont plus décalés d'un même intervalle comme en 7 : le lever
    est avancé de la différence des longitudes augmentée de celle des
    différences ascensionnelles, le coucher de la différence des longitudes
    diminuée de celle-ci.\footnote{Le feuillet dit seulement qu'il s'ensuit
    « l'inégalité des arcs, mais non l'anticipation des levers et des couchers »
    (« Non autem ortuum et occasuum anticipatio ») ; pris à la lettre ce serait
    faux, puisque les levers sont en général décalés. On donne l'énoncé exact.
    « inclinatũ » est écrit avec un trait sur le u (vue 12, gauche).}
  \item[10.] Au pôle, le Soleil reste six mois sur l'horizon et six mois
    au-dessous, si l'on néglige, comme le feuillet le dit expressément,
    l'inégalité de son mouvement.
  \item[11.] Du pôle au cercle polaire, ce jour continu décroît de six mois à
    vingt-quatre heures : autant de mois que de signes, autant de jours que de
    degrés dans l'arc d'écliptique qui ne se couche pas.
  \item[12.] Sur le cercle polaire, le plus long jour est de vingt-quatre heures
    avec un instant de nuit, et la plus longue nuit de même. En tout lieu, le
    temps total de jour sur l'année est d'une demi-année — continu au pôle,
    fractionné ailleurs.
\end{enumerate}

La dernière affirmation est vraie dans le modèle que pose la proposition 10 :
Soleil ponctuel, sans réfraction, parcourant l'écliptique d'un mouvement
uniforme. Deux points diamétralement opposés de l'écliptique ont des
déclinations opposées, donc l'arc diurne de l'un égal à l'arc nocturne de
l'autre, et le Soleil passe autant de temps sur l'un que sur l'autre. Avec le
mouvement réel du Soleil, plus lent en été boréal, l'hémisphère nord reçoit un
peu plus d'heures de jour que l'hémisphère sud.\footnote{Le feuillet s'appuie sur
« 23 secundi Sphæricorum Theodosii » pour l'égalité des arcs « coalternes » ;
l'hypothèse du mouvement uniforme est posée à la proposition 10 (« Supponitur
enim solis motus æqualis, ubi de arcubus primi motus agitur »).}

\section*{Les Phénomènes d'Euclide}

\pagerange{13}{15}
Les \emph{Phénomènes} tirent de l'observation la forme du monde et étudient les
levers et couchers des signes. Le préambule conclut de ce que les étoiles restent
à distance constante et décrivent des parallèles que le monde est une sphère
tournant uniformément autour d'un axe, et définit horizon, méridien, équateur,
zodiaque et tropiques ; l'horizon est un grand cercle, puisqu'il coupe toujours
un grand cercle en deux.\footnote{Une tache d'encre sur « et » (vue 13, gauche).}

\begin{enumerate}
  \item[1.] La Terre est au centre du monde : visant avec la même dioptre deux
    signes opposés, l'un au lever, l'autre au coucher, puis deux autres, on
    obtient deux diamètres du zodiaque qui se coupent au lieu de
    l'observateur.\footnote{« mũdo » est abrégé pour « mundo ».}
  \item[2.] Un cercle passant par les pôles est deux fois par révolution
    perpendiculaire à l'horizon, et l'écliptique deux fois perpendiculaire au
    méridien. L'écliptique n'est jamais perpendiculaire à l'horizon si le zénith
    est hors de la zone comprise entre les tropiques ; elle l'est deux fois par
    jour si le zénith est entre les tropiques, une fois s'il est sur un
    tropique.\footnote{Le feuillet écrit « quoniam polus mundi fuerit extra
    tropicos », là où la démonstration qui suit (« Ibi enim Zodiacus nunquam
    transit per polos horizontis hoc est per verticem loci ») demande le pôle de
    l'horizon, c'est-à-dire le zénith ; et « Si vero polus horizontis in tropico
    fuerit, Zodiacus circulus ad horizontem bis rectus erit », alors que la
    preuve, où le parallèle du zénith coupe l'écliptique en deux points, ne vaut
    que pour un zénith strictement entre les tropiques. Sur le tropique, ce
    parallèle est tangent et l'écliptique n'est perpendiculaire à l'horizon
    qu'une fois par jour, comme le dit la proposition 3 des
    \emph{Habitations}.}
  \item[3.] Une étoile fixe se lève et se couche toujours aux mêmes points de
    l'horizon.
\end{enumerate}

Sous le titre « Additio duarum propositionum », Maurolico en ajoute six, qui
sont celles de l'horizon oblique :\footnote{Le titre annonce deux propositions ;
la page en compte six, numérotées 4 à 9 (vue 13, droite).}

\begin{enumerate}
  \item[4.] Des astres situés sur un même cercle passant par les pôles se lèvent
    et se couchent ensemble sur l'horizon droit.
  \item[5.] Sur un horizon oblique, les astres situés sur la moitié orientale
    d'un grand cercle tangent au plus grand des parallèles toujours visibles se
    lèvent ensemble ; ceux de la moitié occidentale se couchent ensemble.
  \item[6.] Pour des astres situés sur un grand cercle qui ne touche ni ne coupe
    ce parallèle, celui qui se lève le premier se couche le premier.\footnote{Le
    feuillet répète la même moitié d'énoncé (« quæ prius oriuntur, prius
    occidunt : et quæ prius oriuntur prius occidunt ») ; la démonstration
    établit l'énoncé donné ici. « occidentalibus » est écrit là où la
    construction attend un singulier.}
  \item[7.] Pour des astres situés sur un grand cercle qui coupe ce parallèle,
    celui qui est le plus au nord se lève le premier et se couche le dernier —
    « Euclide parle pour notre situation » ; chez les antèques, c'est l'astre le
    plus proche de leur pôle.
  \item[8.] Deux astres diamétralement opposés sur un grand cercle se lèvent
    quand l'autre se couche.
  \item[9.] Le zodiaque se lève sur toute la partie de l'horizon comprise entre
    les tropiques ; et si le plus grand parallèle toujours visible n'est pas plus
    petit que le tropique — latitude au moins égale à celle du cercle polaire —,
    il se lève sur toute la moitié orientale de l'horizon, qui est alors tout
    entière entre les tropiques.\footnote{« minor » est surchargé d'encre et
    « orientalem » couvert d'une tache au milieu (vue 14, gauche). La lecture
    « non minor » est celle qui rend la démonstration juste : l'horizon a pour
    déclinaisons extrêmes $\pm(90^\circ-\varphi)$, comprises entre les
    tropiques dès que $\varphi \ge 90^\circ-\varepsilon$.}
\end{enumerate}

Les propositions d'Euclide reprennent ensuite :

\begin{enumerate}
  \item[10.] Des signes égaux se lèvent et se couchent sur des segments
    d'horizon inégaux : les plus grands près de l'équateur, les plus petits
    près des tropiques ; égaux pour des signes également distants de
    l'équateur.\footnote{Le renvoi est écrit « Hoc autem ostendit Theodosius in
    propositionibus 3. 5. 7 et 9. tertii Sphæricorum Menelai aut in 46 2. » :
    Théodose et Ménélas sont confondus, et « 46 2 » semble renvoyer à la
    proposition 46 d'un livre II.}
  \item[11.] Les demi-cercles de l'écliptique qui commencent en des parallèles
    différents se lèvent en des temps inégaux, le plus long étant celui qui
    commence avec le Cancer, le plus court celui qui commence avec le
    Capricorne (pour l'hémisphère nord) ; deux demi-cercles commençant sur un
    même parallèle se lèvent en temps égaux. Le temps de lever d'un demi-cercle
    est en effet l'arc diurne de son point de départ.
  \item[12.] Si deux demi-cercles ayant un arc commun se lèvent en des temps
    différents, les arcs restants se lèvent en des temps qui diffèrent d'autant ;
    s'ils se lèvent en temps égaux, les restes aussi.\footnote{« subtrahitur etiam
    eius tempus » : « eius » rend un mot abrégé surmonté d'un trait, qui se lit
    à peu près « eoc » (vue 14, droite).}
  \item[13.] Deux arcs égaux et opposés de l'écliptique : pendant que l'un se
    lève, l'autre se couche.\footnote{« reliqua \emph{occidit} oritur » : le
    premier verbe biffé.}
\end{enumerate}

Puis « Additio trium propositionum » :

\begin{enumerate}
  \item[14.] Pour deux horizons « semblables » (tous deux droits, ou de même
    latitude), un astre se lève, culmine et se couche plus tôt sur l'horizon le
    plus oriental, du même intervalle : trente degrés d'équateur font deux
    heures.
  \item[15.] Les arcs d'équateur compris entre les moitiés orientales de tels
    horizons se lèvent avec les arcs d'écliptique compris entre les mêmes
    moitiés.
  \item[16.] Sur l'horizon droit, des arcs égaux de l'écliptique se lèvent (et
    se couchent) en un temps d'autant plus long qu'ils sont plus près des
    solstices, et en temps égaux s'ils sont également distants d'un équinoxe.
    \emph{Corollaire} : dans la sphère droite, les Gémeaux, le Cancer, le
    Sagittaire et le Capricorne se lèvent dans les temps les plus longs, le
    Taureau, le Lion, le Scorpion et le Verseau dans des temps moindres, les
    Poissons, le Bélier, la Vierge et la Balance dans les plus courts.
  \item[17.] Sur un horizon oblique boréal, parmi les arcs égaux de la moitié de
    l'écliptique qui commence au Cancer, celui qui touche le tropique se couche
    dans le temps le plus long, puis les suivants dans des temps décroissants
    jusqu'à l'équinoxe ; des arcs également distants de l'équinoxe se couchent
    en temps égaux. \emph{Corollaire} : sur notre horizon, le Cancer et le
    Sagittaire se couchent dans les temps les plus longs, le Lion et le Scorpion
    dans des temps moindres, la Vierge et la Balance dans les plus courts.
  \item[18.] Parmi les arcs égaux de la moitié qui commence au Capricorne, celui
    qui touche le tropique se lève dans le temps le plus long, jusqu'à
    l'équinoxe dans des temps décroissants. \emph{Corollaire} : le Capricorne et
    les Gémeaux se lèvent dans les temps les plus longs, le Verseau et le
    Taureau dans des temps moindres, les Poissons et le Bélier dans les plus
    courts.
\end{enumerate}

Les trois corollaires sont exacts. On l'a vérifié en calculant les ascensions
obliques par signe à la latitude de Messine ($38^\circ$, obliquité
$23^\circ 30'$) : les signes se lèvent en temps égaux deux à deux, symétriquement
par rapport à l'équinoxe (Bélier et Poissons $18{,}7^\circ$, Taureau et Verseau
$22{,}4^\circ$, Gémeaux et Capricorne $29{,}0^\circ$, Cancer et Sagittaire
$35{,}4^\circ$), et un signe se couche dans le temps où se lève le signe
opposé.\footnote{Le temps de lever est compté ici en degrés d'équateur. Au
corollaire de la proposition 18, la transcription remarque que « Gemini » et
« Taurus » sont écrits là où le parallèle avec le corollaire précédent
attendrait « Sagittarius » et « Scorpius ». Le feuillet a raison : pour les
\emph{levers}, les Gémeaux vont avec le Capricorne et le Taureau avec le Verseau ;
c'est pour les \emph{couchers} que le Sagittaire va avec le Cancer.} Ils ne
comparent que les signes d'une même moitié : à $38^\circ$, le Lion et le Scorpion
($37{,}5^\circ$) se lèvent un peu plus lentement que la Vierge et la Balance
($37{,}1^\circ$), de sorte qu'un ordre énoncé sur le zodiaque entier serait
faux ; le feuillet ne l'énonce pas.\footnote{Pour les signes « au-delà de
l'équateur », le feuillet dit de substituer le Capricorne, le Verseau et les
Poissons au Cancer, au Lion et à la Vierge. Dans la démonstration de la
proposition 17, « et occidunt » à la fin de l'énoncé est écrit ainsi ; les
renvois sont « per 6. et 10 tertij Sphæricorum Theod. et 46. Sphæricorum
Menelai ». « horizontib\textsuperscript{9} » abrège « horizontibus » (vue 15,
droite).}

Maurolico conclut en une phrase qui dit sa méthode : il a démontré ainsi les
\emph{Phénomènes} d'Euclide, avec l'aide des \emph{Sphériques} de Ménélas et de
Théodose, « beaucoup plus facilement que l'auteur ne l'avait fait », d'où il
paraît combien les \emph{Sphériques} sont non seulement commodes, mais
nécessaires pour démontrer les rudiments de l'astronomie.\footnote{« Sic
Euclidis Phoenomena et Menelai et Sphæricorum Theodosii adiumentis ubi multo
facilius quam autor fecerat demonstravimus, unde patet quanta sit non solum
commoditas sed etiam necessitas Sphæricorum in astronomicis rudimentis
demonstrandis » (vue 15, droite).}

\section*{Les habitants de la Terre, et la visibilité des astres}

\pagerange{15}{16}
« Habitationum collatio » classe les habitants d'après ce qu'ils voient du ciel.
Les périèques, sous le même parallèle, ont même hauteur du pôle, mêmes jours et
mêmes saisons ; les antèques, sous des parallèles égaux et opposés, ont les
saisons inversées et des ombres méridiennes égales mais dirigées en sens
contraires ; les antipodes, en outre diamétralement opposés, ont le même horizon
et voient se lever ce qui se couche pour les autres. Les amphisciens, entre les
tropiques, ont l'ombre de midi tantôt au nord tantôt au sud, et « deux étés et
deux hivers » ; les périsciens, au pôle, voient l'ombre tourner tout autour
d'eux.\footnote{Le feuillet donne ensuite « Antiscij sunt perioeci quibus quidem
æqualis et eodem versus meridiana umbra flectitur » et « Heteroscij sunt qui et
Antoeci quibus in diversum proijcitur umbra meridiana », puis un second sens de
« hétérosciens » : ceux dont l'ombre de midi tombe toujours du même côté, hors des
tropiques. C'est ce second sens qui est l'usage courant ; « antisciens » désigne
d'ordinaire les habitants de part et d'autre de l'équateur dont les ombres de
midi sont opposées, ce que le feuillet appelle ici hétérosciens (vue 16,
gauche).}

« De Astrorum fulsionibus » corrige la notion d'intervalle de visibilité : la
distance qui sépare du Soleil un astre au moment de sa première ou de sa
dernière apparition ne se mesure pas sur le zodiaque, mais sur le cercle vertical,
comme abaissement du Soleil sous l'horizon, et elle n'est pas la même pour tous
les astres, contrairement à ce qu'Autolycus a supposé : plus l'astre est
lumineux, plus elle est petite. Ceux qui disent la Lune invisible à moins de
14° du Soleil se trompent ; Pline a tort d'en conclure que la Lune est plus
petite que les autres astres, car un grand intervalle dénote un éclat apparent
faible, non une petite taille. Tout cela est juste.

Il en tire un calcul qui ne l'est pas. Vénus, selon Ptolémée, demande
$5^\circ$. Si la Lune demande $x$ degrés, elle aurait alors $x/180$ de la
lumière de la pleine Lune ; la pleine Lune valant cent fois Vénus (le rapport
des diamètres étant dix, selon Albategnius), le croissant vaudrait
$100x/180$ fois Vénus, et les intervalles étant inverses des lumières,
$x = 5 \cdot 180/(100x)$, d'où $x^2 = 9$ et $x = 3^\circ$. Le feuillet présente
ce raisonnement en partant de $3^\circ$ et en retrouvant $3^\circ$, ce qui en fait
une vérification de cohérence, dont $3^\circ$ est bien la seule solution
positive.\footnote{« Nam cum tantum interstitium sit sexagesima pars semicirculi,
ut tunc luna adepta sit de plenilunio toti lumini sexagesimam. Sed plena Luna
centupla ferme est venereæ faciei […] Et ideo illud lunæ corniculatum lumen tenue
ad lumen veneris erit sicut quinque ad 3 » (vues 16, gauche et droite). « toti »
est traversé d'un trait oblique, peut-être une biffure ; « cuius » biffé avant
« illius ». L'arithmétique est juste : $100/60 = 5/3$ et
$5 \times 3/5 = 3$.} Mais aucune prémisse ne tient. La fraction éclairée du
disque lunaire à l'élongation $\psi$ est $(1-\cos\psi)/2$, non $\psi/180^\circ$ :
à $3^\circ$ elle vaut moins d'un millième, et non un soixantième ; l'abaissement du
Soleil sous l'horizon n'est pas l'élongation ; le seuil de visibilité n'est pas
inversement proportionnel à l'éclat. Le croissant n'est en fait jamais vu à
moins de sept degrés environ du Soleil. La conséquence que le feuillet en tire —
une Lune visible le matin et le soir d'un même jour si la conjonction tombe vers
midi en été — n'est donc pas établie.

La table des sinus, qui occupe le reste de la page, ouvre la section suivante.

\section*{Les tables : sinus, tangentes, sécantes, et la table des trois arcs}

\pagerange{16}{20}
Quatre tables occupent les vues 16 à 20. Les trois premières sont, sous leurs
noms du XVI\textsuperscript{e} siècle, nos tables de sinus, de tangentes et de
sécantes, de degré en degré de $1^\circ$ à $90^\circ$, pour un rayon (« sinus
totus ») de $100\,000$. On écrit dans toute la suite $\sin$, $\cos$, $\tan$,
$\sec$ pour les lignes rapportées au rayon 1 ; une entrée de table est donc
$100\,000$ fois la valeur moderne, arrondie.\footnote{Le « sinus rectus » d'un arc
est la demi-corde de l'arc double, notre $R\sin$ ; le « sinus secundus » est le
sinus du complément, notre $R\cos$ ; le « sinus versus » est la flèche,
$R(1-\cos)$, qu'on note $\operatorname{vers}$. Le rayon $R = 100\,000$ est ce que
les textes appellent « sinus totus » ; diviser par lui revient, disent-ils, à
« rejeter cinq figures à droite ». Ces équivalences sont posées une fois pour
toutes et ne sont plus rappelées.} Chaque table a quatre colonnes : le degré,
la valeur, la différence avec la ligne précédente, et le soixantième de cette
différence, qui est l'accroissement par minute d'arc et sert à interpoler. Ainsi
$\sin 23^\circ 30'$ se tire de $\sin 23^\circ = 39\,073$ et du soixantième
$26{:}41$ de la ligne suivante : $39\,073 + 30 \times 26\tfrac{41}{60} \approx
39\,874$, et les problèmes emploient $39\,875$.

\subsection*{La table des sinus}

\pagerange{16}{16}
« Tabella sinus recti » donne $100\,000\sin n^\circ$. Recalculée, elle est exacte
à une unité près sur chaque ligne ; l'écart d'une unité, aux lignes 48, 67, 73 et
85, est affaire d'arrondi. Trois erreurs touchent les colonnes auxiliaires :

\begin{itemize}
  \item ligne 59, le soixantième de la différence 912 est écrit $15{:}21$ ; il
    vaut $15{:}12$ ;
  \item ligne 86, la différence est écrite 139 ; $99\,756 - 99\,620 = 136$, et le
    soixantième écrit, $2{:}16$, est bien celui de 136 ;
  \item ligne 88, le soixantième de 76 est écrit $1{:}6$ ; il vaut
    $1{:}16$.\footnote{La table est en deux blocs de quatre colonnes, intitulées
    au pied « gradus », « sinus », « Dria » (abrégé pour « Differentia »),
    « 6\textsuperscript{ma} » et, à droite, « 60\textsuperscript{ma} ». Aux lignes
    46 et 47, un point au lieu des deux-points. Chiffres repassés ou surchargés,
    dont la transcription donne la lecture finale : lignes 10, 24, 35, 39 (le 6
    ajouté au-dessus), 48, 59, 63. Ligne 4, le troisième chiffre de 6976 est formé
    presque comme un 1 ; la lecture 4 est appuyée par la colonne suivante et par
    la valeur exacte. La ligne 90 est écrite « 100.000 » (vue 16, droite).}
\end{itemize}

\subsection*{La « table féconde » : les tangentes}

\pagerange{17}{17}
« Tabula foecunda » donne, sous le nom d'« umbra », l'ombre d'un gnomon égal au
rayon : $100\,000\tan n^\circ$.\footnote{Le mot est défini plus loin dans le
volume (vue 44, et proposition XXX du second livre, vue 46) : pour un gnomon
égal au rayon, l'« umbra versa » d'un arc est notre tangente, l'« umbra recta »
notre cotangente. La « table féconde » est dite, à la vue 44, celle « de
Iohannes de Monteregio », Regiomontanus ; le feuillet de la vue 17 ne le dit
pas.} Les valeurs sont exactes à trois unités près, sauf :

\begin{itemize}
  \item ligne 40, le numéro du degré est écrit « 46 » ;
  \item ligne 43, la valeur est écrite $63\,252$ ; c'est $93\,252$, comme le
    veulent la différence 3212 et la valeur exacte ;
  \item ligne 50, $119\,197$ au lieu de $119\,175$ : l'erreur, de 22 unités, est
    dans la valeur elle-même, puisque les deux différences voisines ($4160$ et
    $4294$ au lieu de $4138$ et $4316$) en tiennent compte ;
  \item ligne 78, $470\,453$ au lieu de $470\,463$, de même avec sa
    différence ;
  \item ligne 84, la différence est écrite $137\,061$ ; elle vaut $137\,001$,
    dont $2283{:}21$ est bien le soixantième ;
  \item soixantièmes : ligne 34, $41{:}51$ pour $41{:}52$ ; ligne 68,
    $196{:}44$ pour $198{:}44$ ; ligne 69, « $216{:}4$ » pour $216{:}40$ ;
    ligne 83, $1714{:}18$ pour $1714{:}58$.\footnote{Colonnes intitulées en tête
    « gr. », « umbra », « Dria », « 6\textsuperscript{ma} ». Ligne 69, la
    transcription lit « 216: 4 » sans chiffre après le 4. Ligne 4, dernière
    colonne, « 11 » est une lecture douteuse. Chiffres surchargés ou repassés
    aux lignes 16, 38, 68 (numéro du degré), 74, 75. Le point ou les deux-points
    de la dernière colonne sont gardés tels qu'écrits dans la transcription ;
    ils ne changent pas la valeur (vue 17, gauche).}
\end{itemize}

Sous la table, cinq valeurs pour les arcs voisins de $90^\circ$, où
l'interpolation n'est plus possible : $\tan 89^\circ 15' = 7\,638\,998$,
$\tan 89^\circ 30' = 11\,458\,872$, $\tan 89^\circ 45' = 22\,918\,163$,
$\tan 89^\circ 55' = 68\,754\,439$, $\tan 89^\circ 59' = 343\,772\,546$. Les
trois premières sont justes aux dernières unités près ; les deux dernières sont
un peu faibles (les valeurs exactes sont $68\,754\,887$ et $343\,774\,667$), avec
une erreur relative de l'ordre de six millionièmes.

\subsection*{La « table bienfaisante » : les sécantes}

\pagerange{19}{19}
« Tabella benefica » donne, sous le nom de « Radius », le rayon mené du pied du
gnomon à l'extrémité de l'ombre, c'est-à-dire $100\,000\sec n^\circ$ ; Maurolico
dit l'avoir faite « à l'imitation de la table féconde ».\footnote{Vue 44, gauche,
« Demonstratio tabulæ beneficæ » : voir la section sur Regiomontanus. Colonnes
intitulées en tête « g. », « Radius », « Dria », « 60\textsuperscript{ma} » ; au
second bloc, la ligne 90 porte « Infinitum » dans les trois colonnes. La
transcription ne reproduit pas les points qui précèdent la plupart des
différences ni ceux qui séparent parfois les milliers (vue 19, droite, numéro à
l'encre 18).} Les valeurs sont exactes à une unité près, sauf :

\begin{itemize}
  \item ligne 16, $104\,630$ ; la valeur est $104\,030$, que donne aussi la
    différence 502 ajoutée à la ligne précédente ;
  \item ligne 24, $108\,464$, chiffre surchargé ; la valeur est $109\,464$ ;
  \item ligne 30, la différence est écrite 1130 ; $115\,470 - 114\,335 = 1135$,
    et le soixantième écrit, $18.55$, est celui de 1135, non de 1130 ;
  \item ligne 9, différence 264 pour 263 ;
  \item ligne 4, soixantième « 2.47 » pour $1{:}47$ ;
  \item dans la dernière colonne du second bloc, les lignes 51 à 55 répètent les
    valeurs des lignes 46 à 50, et de la ligne 56 à la ligne 85 chaque valeur est
    le soixantième de la différence écrite cinq lignes plus haut : la colonne a
    glissé de cinq lignes à la copie, et se remet d'accord aux lignes 86 à
    89.\footnote{Constats de la transcription, qui a contrôlé chaque rayon par
    l'addition des différences ; les valeurs exactes sont de cette lecture.
    Ligne 25, le 3 de $110\,338$ est surchargé ; ligne 44, un trait de trop après
    $139\,016$ ; ligne 43 de la dernière colonne, « 36.10 » se lit aussi
    « 39.20 » ; ligne 86, « 4769.47 », le quatrième chiffre se lit aussi 4 ;
    ligne 89, « 47741.37 », le cinquième chiffre, surchargé, se lit aussi 5.}
\end{itemize}

Sous la table, six valeurs : $\sec 0^\circ = 100\,000$, $\sec 89^\circ 15' =
7\,639\,653$, $\sec 89^\circ 30' = 11\,459\,309$, $\sec 89^\circ 45' =
22\,918\,381$, $\sec 89^\circ 55' = 68\,754\,512$, $\sec 89^\circ 59' =
343\,772\,560$ ; comme pour les tangentes, les deux dernières sont faibles de
quelques millionièmes (valeurs exactes $68\,754\,960$ et $343\,774\,682$).

\subsection*{La table des trois arcs}

\pagerange{20}{20}
La page de gauche de la vue 20 porte une table sans titre donnant, pour chaque
degré $\lambda$ de longitude écliptique, trois arcs en degrés et minutes,
« Primus », « Secundus », « Tertius ». Les problèmes 4 et 5 disent employer une
« table générale des déclinaisons et des ascensions » à trois arcs et les
définissent ; c'est celle-ci.\footnote{Le rapprochement est fait par la
transcription (vue 20, gauche), à partir de l'usage des trois arcs dans les
problèmes 4 et 5.} Menons par le point $\lambda$ de l'écliptique le cercle de
latitude, perpendiculaire à l'écliptique, jusqu'à l'équateur. Le triangle
sphérique formé avec l'équinoxe est rectangle au point de l'écliptique, son
angle à l'équinoxe est l'obliquité $\varepsilon$, et ses trois arcs sont :

\begin{itemize}
  \item \emph{Primus}, l'arc $p$ du cercle de latitude entre l'écliptique et
    l'équateur : $\tan p = \tan\varepsilon \sin\lambda$ ;
  \item \emph{Secundus}, l'angle $\theta$ de ce cercle avec l'équateur :
    $\cos\theta = \sin\varepsilon\cos\lambda$ ;
  \item \emph{Tertius}, l'arc $t$ d'équateur entre l'équinoxe et ce cercle :
    $\tan t = \tan\lambda / \cos\varepsilon$.
\end{itemize}

Avec $\varepsilon = 23^\circ 30'$, que les problèmes emploient, les quatre-vingt-dix
lignes s'accordent avec ces formules à moins de deux minutes près, à cinq
exceptions : ligne 5, \emph{Primus} $2^\circ 20'$ pour $2^\circ 10'$ ; ligne 37,
\emph{Primus} $14^\circ 30'$ pour $14^\circ 40'$ ; ligne 68, \emph{Secundus}
$81^\circ 22'$ pour $81^\circ 25'$ ; ligne 70, \emph{Secundus} $82^\circ 20'$ pour
$82^\circ 10'$ ; ligne 89, \emph{Secundus} $89^\circ 33'$ pour
$89^\circ 36'$.\footnote{Ligne 5, la transcription note que le 2 des minutes est
précédé d'un point et se lit peut-être 1 : la valeur exacte appuie « 10 ». Lignes
17 à 20, colonne \emph{Secundus}, les minutes sont surchargées de deux tracés à
la fois (35 et 38, 43 et 48, 51 et 54, 50 et 59) ; la transcription offre 35, 43,
51, 59 en suivant la progression de la colonne, et les valeurs exactes
($67^\circ 35{,}1'$, $42{,}8'$, $51{,}0'$, $59{,}6'$) confirment ce choix. Ligne 79,
\emph{Tertius}, « 79. 54 », dont le dernier chiffre se lit aussi 9 : la valeur
exacte est $79^\circ 53{,}6'$. Ligne 18, « 7. 39 », le dernier chiffre se lit aussi
4 : la valeur exacte est $7^\circ 39{,}2'$. Chiffres surchargés mais lisibles aux
lignes 16, 27, 43, 53, 63, 83.} À $\lambda = 90^\circ$ la table donne
$23^\circ 30'$, $90^\circ$, $90^\circ$ ; au début du Bélier, le \emph{Secundus}
vaut $90^\circ - \varepsilon = 66^\circ 30'$, comme le rappelle le problème 5, et la
ligne 1 porte déjà cette valeur à la minute près.

\section*{Huit problèmes résolus avec les tables}

\pagerange{17}{19}
Les problèmes sont numérotés en marge. Chacun donne une règle de calcul, puis un
exemple entièrement chiffré. On donne ici chaque règle sous sa forme moderne, en
disant quel triangle sphérique la justifie — le feuillet ne le dit pas, le second
livre des \emph{Sphériques} le fait plus loin —, puis on refait
l'exemple.\footnote{Les problèmes 1 à 6 se lisent dans l'ordre du texte sur la
vue 18, droite (numéro à l'encre 17 ; problèmes 1 à 4), la vue 19, gauche (fin
du 4, et 5), la vue 17, droite (numéro 16 ; fin du 5), et la vue 18, gauche
(problème 6) : le problème 5 s'arrête au bas de la vue 19 sur « Nam in primis »
et reprend en tête de la vue 17, droite, sur « considerandum est ». Les deux
feuillets numérotés 16 et 17 sont reliés dans l'ordre inverse de celui du texte ;
la reconstitution est celle de la transcription, qui garde l'ordre des vues.}
Dans toute la suite, $\lambda$ est la longitude écliptique, $\beta$ la latitude
écliptique, $\delta$ la déclinaison, $\alpha$ l'ascension droite, $\varphi$ la
latitude du lieu, $\varepsilon = 23^\circ 30'$ l'obliquité, et « rejeter cinq
figures » veut dire diviser par $100\,000$.

\subsection*{Problème 1 : la déclinaison d'un point de l'écliptique}

\emph{Règle.} $\sin\delta = \sin\lambda\,\sin\varepsilon$ — le triangle rectangle
formé par l'équateur, l'écliptique et le cercle de déclinaison du point.
\emph{Exemple}, la fin du onzième degré du Taureau, $\lambda = 41^\circ$ :
$65\,606 \times 39\,875 = 2\,616\,039\,250$, soit $26\,160$, et $\delta =
15^\circ 10'$. C'est exact.\footnote{« Taurus » est écrit pour « Tauri » ; le 15 de
« 15 graduum » est en partie recouvert par le g de « graduum » ; le problème
suivant le reprend en « g. 15. m. 10. » (vue 18, droite).}

\subsection*{Problème 2 : l'ascension droite, par la table féconde}

\emph{Règle.} $\sin\alpha = \tan\delta\,\tan(90^\circ-\varepsilon) =
\tan\delta\cot\varepsilon$, l'arc étant compté depuis l'équinoxe le plus proche.
Le feuillet ajoute que, de même, la déclinaison d'un astre et la latitude du lieu
donnent par la table féconde sa différence ascensionnelle :
$\sin\Delta = \tan\delta\,\tan\varphi$. \emph{Exemple} : $\tan 15^\circ 10' =
27\,108$ et $\tan 66^\circ 30'$, écrit $230\,094$ ; le produit $6\,237\,388\,152$
donne $62\,374$, et le feuillet conclut à $\alpha = 38^\circ$.

La tangente de $66^\circ 30'$ vaut $229\,984$, non $230\,094$. Avec la bonne
valeur, le produit est $6\,234\,406\,272$, le sinus $62\,344$, et
$\alpha = 38^\circ 34'$, que le problème 3 retrouve par une autre voie ; la valeur
exacte est $38^\circ 33{,}7'$.\footnote{La transcription lit « 230094 » en
signalant que l'avant-dernier chiffre est mal formé, et remarque que le produit
écrit est bien celui de $27\,108$ par $230\,094$ : l'erreur porte donc sur la
valeur prise dans la table, non sur la lecture. Avec $230\,094$, l'arc serait
$38^\circ 35'$ ; le feuillet ne donne que les degrés.}

\subsection*{Problème 3 : l'ascension droite, par la table bienfaisante}

\emph{Règle.} $\cos\alpha = \sec\delta\,\cos\lambda$ : on multiplie la sécante de
la déclinaison par le sinus du complément de l'arc d'écliptique, et l'arc dont le
résultat est le sinus, ôté du quadrant, est l'ascension droite. C'est
$\cos\lambda = \cos\delta\cos\alpha$ dans le même triangle rectangle.
\emph{Exemple} : $\sec 15^\circ 10' = 103\,612$ (valeur interpolée ; exacte,
$103\,609$) et $\cos 41^\circ = 75\,471$ ; le produit $7\,819\,701\,252$ donne
$78\,197$, arc $51^\circ 26'$, complément $38^\circ 34'$. C'est
juste.\footnote{Le feuillet écrit « 75471 sinum scilicet graduum 41 », là où il
s'agit du sinus du complément de $41^\circ$, c'est-à-dire du sinus de $49^\circ$
que donne la table ; la règle, elle, dit bien « sinum secundum ».}

\subsection*{Problème 4 : la déclinaison d'une étoile}

\emph{Règle.} On entre dans la table des trois arcs avec la longitude de l'étoile,
comptée depuis l'équinoxe le plus proche. L'arc $p$ (\emph{Primus}), de même nom
que l'hémisphère où est le point de l'écliptique, se combine avec la latitude
$\beta$ de l'étoile : somme si elles sont de même nom, différence sinon, avec le
nom de la plus grande. Le résultat $s = p \pm \beta$, que le feuillet appelle
« argument de la déclinaison », est l'arc du cercle de latitude entre l'équateur
et l'étoile ; ce cercle coupant l'équateur sous l'angle $\theta$
(\emph{Secundus}), la déclinaison est donnée par
\[
\sin\delta = \sin s\,\sin\theta ,
\]
dans le triangle rectangle formé par ce cercle, l'équateur et le cercle de
déclinaison de l'étoile ; elle a le nom de l'argument. Le sinus de $\theta$ est,
dit le feuillet, « le nombre multiplicande dont se sert Regiomontanus ».

\emph{Exemple}, Aldebaran, « l'œil du Taureau », dont la longitude est « en ce
temps » de $2^\circ 50'$ des Gémeaux ($\lambda = 62^\circ 50'$) et la latitude
$5^\circ 10'$ sud. La table, interpolée, donne $p = 21^\circ 9'$,
$\theta = 79^\circ 31'$, $t = 64^\circ 47'$ ; les formules donnent
$21^\circ 8{,}9'$, $79^\circ 30{,}6'$, $64^\circ 47{,}8'$. L'argument est
$21^\circ 9' - 5^\circ 10' = 15^\circ 59'$ nord, $\sin s = 27\,536$,
$\sin\theta = 98\,331$ ; le produit, $2\,707\,642\,416$, donne $27\,076$, et
$\delta = 15^\circ 43'$ nord, à une demi-minute de la valeur exacte pour ces
données.\footnote{Le feuillet écrit le produit « 2707642916 » ; la division par
$100\,000$ n'en est pas affectée. Il écrit d'abord « grad. 21. m. 5. » et
« g. 79. m. 10. », puis emploie $21^\circ 9'$ et $79^\circ 31'$, qui sont les
bonnes interpolations. Le mot rendu « vel », abrégé « ul », se lit mal ;
« gratia » manque après « Exempli » ; « cum ambobus » est écrit deux fois, la
première biffée ; « septen\textsuperscript{lis} » abrège « septentrionalis »
(vue 19, gauche).}

\subsection*{Problème 5 : l'ascension droite d'une étoile}

\emph{Règle.} Soit $\Delta$ l'arc d'équateur compris entre le cercle de latitude
et le cercle de déclinaison de l'étoile, que le feuillet nomme « différence du
passage de l'étoile au méridien ». Dans le même triangle rectangle,
\[
\sin\Delta = \tan\delta\,\cot\theta \qquad\text{ou}\qquad
\cos\Delta = \frac{\cos s}{\cos\delta} = \sec\delta\,\cos s ,
\]
la première par la table féconde (avec $\delta$ et $90^\circ-\theta$), la seconde
par la table bienfaisante. L'ascension droite s'obtient alors en combinant $\Delta$
avec l'arc $t$ (\emph{Tertius}) selon le quadrant de l'écliptique où tombe
l'étoile et le nom de sa déclinaison.\footnote{Dans la table des trois arcs,
« ingredere tabulam declinationum et ascensionum generalem » ; la règle de
combinaison occupe toute la vue 17, droite, et le début de la vue 18, gauche. Deux
lignes y sont biffées (« Si arcus tertius minor quam differentia, tunc arcus
relictus de toto arcus relictus iunctus ») ; « relinquet » est écrit au-dessus de
« conficiet » biffé ; « At » et le second « relinquet » sont des lectures
incertaines ; « dria » et « declinaẽ » abrègent « differentia » et
« declinatione ».}

Les huit cas se ramènent à une seule observation : le cercle de latitude monte
de l'équateur vers le pôle nord de l'écliptique, qui est à l'ascension droite
$270^\circ$, et vers le pôle sud, à $90^\circ$. Pour une étoile boréale des deux
premiers quadrants, ou australe des deux derniers, $\Delta$ se retranche de $t$ ;
dans les autres cas, il s'y ajoute ; puis l'arc $t \mp \Delta$, compté depuis
l'équinoxe le plus proche, donne $\alpha$ selon le quadrant ($t\mp\Delta$,
$180^\circ - (t\mp\Delta)$, $180^\circ + (t\mp\Delta)$, $360^\circ - (t\mp\Delta)$),
avec passage de l'autre côté de l'équinoxe quand $\Delta > t$. On a vérifié les
huit prescriptions du feuillet sur ce principe : toutes sont
justes.

Au début du Bélier ou de la Balance, $p$ et $t$ sont nuls, $\theta =
90^\circ - \varepsilon = 66^\circ 30'$, et l'argument est la latitude $\beta$. Au
Bélier, $\alpha = \Delta$ pour une étoile australe, $360^\circ - \Delta$ pour une
boréale, comme le dit le feuillet. À la Balance, $\alpha = 180^\circ + \Delta$ pour
une étoile boréale et $180^\circ - \Delta$ pour une australe.\footnote{Le
feuillet écrit d'abord « in principio Libræ cum Arietis », où le sens et la suite
demandent le seul Bélier ; pour la Balance boréale, il écrit « arcus sæpe
memoratus cum semicirculo eamdem residuabit », où « residuabit », qu'il emploie
ailleurs pour un reste, ne dit pas l'addition qu'il faut ; le cas de la Balance
australe n'est pas écrit (vue 18, gauche).} Aux solstices, cercle de latitude et
cercle de déclinaison sont le colure solsticial : $\delta = \varepsilon \pm \beta$
et $\alpha = 90^\circ$ ou $270^\circ$, pourvu que $\varepsilon + \beta$ reste
inférieur à $90^\circ$.\footnote{Le feuillet dit « dodrans », trois quarts de
cercle, pour $270^\circ$. Deux lignes biffées au début de la vue 18, gauche, de
« Zodiaci » à « ambobus » ; « In reliquo » est écrit deux fois.}

\emph{Exemple}, Aldebaran. Par la table féconde : $\tan 15^\circ 43' = 28\,140$
et $\tan(90^\circ - 79^\circ 31') = \tan 10^\circ 29' = 18\,506$ ; le produit
$520\,758\,840$ donne $5208$, sinus d'un arc « un peu moindre que $3^\circ$ »
($2^\circ 59'$). Par la table bienfaisante : $\sec 15^\circ 43' = 103\,879$ et
$\cos 15^\circ 59' = 96\,134$, produit $9\,986\,303\,786$, soit $99\,863$, sinus de
$87^\circ$, dont le complément est $3^\circ$. L'étoile est dans le premier quadrant
et boréale : $\alpha = 64^\circ 47' - 3^\circ = 61^\circ 47'$. La valeur exacte pour
ces données est $61^\circ 48{,}7'$, l'écart venant de l'arrondi de $\Delta$ à
$3^\circ$.\footnote{Le feuillet écrit le complément de l'arc second « g. 10. m.
27. » ; $90^\circ - 79^\circ 31' = 10^\circ 29'$, et le nombre qu'il prend,
$18\,506$, est bien la tangente de $10^\circ 29'$ ($18\,504$), non celle de
$10^\circ 27'$ ($18\,444$). Il écrit ensuite « præfatam differentiam 3. grad. 3. »
et n'emploie que 3 degrés. « Demo abicio » : le premier mot se lit mal et n'est
pas biffé ; « dicta » est traversé d'un trait ; le 6 de « 9986303786 » est ajouté
au-dessus de la ligne. La valeur $\sec 15^\circ 43'$ exacte est $103\,884$.}

\subsection*{Problème 6 : l'amplitude ortive}

\emph{Règle.} $\sin A = \sin\delta\,\sec\varphi$, où l'amplitude $A$ est l'arc
d'horizon entre le point est et le lever de l'astre : dans le triangle rectangle
formé par l'horizon, l'équateur et le cercle de déclinaison, l'angle à l'équateur
est $90^\circ - \varphi$.\footnote{« latitudo ortiva » : « l'amplitude ortive ».
Dans la règle, « ortiva » est biffé après « latitudine » ; « abrectis » est écrit
pour « abiectis » (vue 18, gauche).} \emph{Exemple}, à Messine
($\varphi = 38^\circ$), pour le Soleil au solstice ($\delta = 23^\circ 30'$) :
$126\,902 \times 39\,875 = 5\,060\,217\,250$, soit $50\,602$, et
$A = 30^\circ 24'$. C'est exact.

\subsection*{Problèmes 7 et 8}

\pagerange{48}{49}
Deux problèmes numérotés « 7 » et « 8 », à la fin du volume, emploient la même
latitude et l'amplitude du problème 6 ; on les lit à sa suite.\footnote{Le
rapprochement est de cette lecture (voir « Ce que le volume contient »). Sur la
vue 48, droite (numéro à l'encre 47), le « 7 » est écrit au-dessus de la
première ligne du paragraphe, le « 8. » dans la marge ; la page s'ouvre sur la
fin d'une phrase, « qui quidem arcus est horizontis inter æquinoctialem, et
solem in tali horizonte orientem », qui définit l'amplitude.}

\emph{Problème 7 : la différence ascensionnelle.} \emph{Règle} : on multiplie la
sécante de la déclinaison par le cosinus de l'amplitude ; l'arc dont le produit
est le sinus, ôté du quadrant, est la différence ascensionnelle $\Delta$. En effet,
dans le triangle rectangle du problème 6, $\cos A = \cos\delta\cos\Delta$. Ajoutée
au quadrant pour une déclinaison boréale, retranchée pour une australe, elle donne
l'arc semi-diurne. \emph{Exemple} : $\sec 23^\circ 30' = 109\,038$ et
$\cos 30^\circ 24' = 86\,251$, produit $9\,404\,636\,538$, soit $94\,046$, arc
$70^\circ 8'$, et $\Delta = 19^\circ 52'$. La valeur exacte, par
$\sin\Delta = \tan\varphi\tan\delta$, est $19^\circ 51{,}6'$.\footnote{Le feuillet
écrit « Ad dicti loci latitudinem gr. 38 ex tabula benefica excipio numerum
109038 » : la phrase annonce la sécante de la latitude, qui vaut $126\,902$, mais
le nombre pris est celui de la déclinaison, comme le veut la règle ; sa valeur
exacte est $109\,044$. La règle dit que l'arc ôté du quadrant « residuabit
ascensionem quæsitam » : c'est la différence ascensionnelle qu'on cherche. Dans
« 940.463.6538 », les chiffres 53 sont surchargés.}

\emph{Problème 8 : la distance au méridien d'un astre dont on connaît la hauteur.}
Le feuillet ne donne pas de règle générale ; il la montre sur un exemple. On
multiplie le sinus de la hauteur $h$ par la sécante de la latitude, puis le
résultat par la sécante de la déclinaison : c'est l'« argument de la distance »,
$\sin h\,\sec\varphi\,\sec\delta$. On le compare au sinus de la différence
ascensionnelle : égal, la distance au méridien est d'un quadrant ; plus petit,
elle dépasse le quadrant de l'arc dont la différence est le sinus ; plus grand,
elle est moindre. Pour une déclinaison australe, on ajoute au lieu de comparer.
C'est exactement la formule moderne de l'angle horaire $H$ :
\[
\cos H = \frac{\sin h - \sin\varphi\sin\delta}{\cos\varphi\cos\delta}
       = \sin h\,\sec\varphi\,\sec\delta \;-\; \tan\varphi\tan\delta ,
\]
où $\tan\varphi\tan\delta = \sin\Delta$ ; le signe du second terme change avec
celui de $\delta$. Si $\delta = 0$, la seconde multiplication est inutile.

\emph{Exemples.} Le Soleil au onzième degré du Taureau, à Messine. Le feuillet
prend $\sec\delta = 103\,609$, $\sec\varphi = 126\,902$, $\sin\Delta = 21\,189$.
Ces trois nombres sont justes pour la déclinaison du problème 1,
$\delta = 15^\circ 10'$ (pour laquelle $\Delta = 12^\circ 14'$). Le feuillet
écrit pourtant « déclinaison $15^\circ 20'$ » et « différence ascensionnelle
$12^\circ 24'$ », deux valeurs qui ne s'accordent ni entre elles ni avec les
nombres employés ; on corrige en $15^\circ 10'$ et $12^\circ 14'$, et les calculs
du feuillet sont alors cohérents.\footnote{Pour $\delta = 15^\circ 20'$, la
sécante serait $103\,691$ et la différence ascensionnelle $12^\circ 22'$ ; le
sinus de $12^\circ 24'$ serait $21\,474$. « Dria » abrège « Differentia » (vue 48,
droite).}

\begin{itemize}
  \item $h = 5^\circ 30'$ : $9585 \times 126\,902 = 1\,216\,355\,670$, soit
    $12\,164$ après arrondi ; $12\,164 \times 103\,609 = 1\,260\,299\,876$, soit
    $12\,603$. Comme $12\,603 < 21\,189$, la différence $8586$ est le sinus de
    $4^\circ 55'$, et $H = 90^\circ + 4^\circ 55' = 94^\circ 55'$. Valeur exacte :
    $94^\circ 55{,}2'$.\footnote{Le nombre intermédiaire est surchargé ; la
    transcription offre $12\,164$ parce que c'est le seul que confirment le
    produit suivant et la règle d'arrondi, ce qui est une vérification et non
    une lecture. L'argument est écrit « 12604 » puis « 12603 » ; $12\,603$ est
    juste. Le résultat est écrit « gr. 95. m. 55 » (vue 49, gauche), et la somme
    annoncée à la fin de la vue 48 donne $94^\circ 55'$.}
  \item $h = 9^\circ 16'$ : $16\,103 \times 126\,902 = 2\,043\,502\,906$, soit
    $20\,435$ ; $20\,435 \times 103\,609 = 2\,117\,249\,915$, soit $21\,172$,
    « presque égal » à $21\,189$ : $H$ est d'un quadrant, à une minute
    près.\footnote{Le feuillet écrit « 2117149915 » et « 21171 ».}
  \item $h = 27^\circ 50'$ : $46\,690 \times 126\,902 = 5\,925\,054\,380$, soit
    $59\,251$ ; $59\,251 \times 103\,609 = 6\,138\,936\,859$, soit $61\,389$ ;
    $61\,389 - 21\,189 = 40\,200 = \cos H$, et $H = 66^\circ 18'$ (exact :
    $66^\circ 17{,}4'$).\footnote{Le feuillet donne le résultat sans l'étape
    intermédiaire : l'arc dont $40\,200$ est le sinus est $23^\circ 42'$, et son
    complément $66^\circ 18'$.}
  \item Même hauteur, déclinaison australe : $61\,389 + 21\,189 = 82\,578$, sinus
    de $55^\circ 40'$, et $H = 34^\circ 20'$ (exact : $34^\circ 20{,}6'$).
\end{itemize}

Le feuillet ajoute trois cas limites, tous justes. Si le parallèle de l'astre
touche seulement l'horizon ($\delta = 90^\circ - \varphi$), on prend le rayon pour
sinus de la différence ascensionnelle, puisque $\tan\varphi\tan\delta = 1$. Si
l'astre est circumpolaire, on retranche du sinus de la hauteur le sinus de sa plus
petite hauteur méridienne $h_0$, et l'on opère sur le reste avec le rayon :
c'est l'identité
\[
\cos H = (\sin h - \sin h_0)\sec\varphi\sec\delta - 1,
\qquad \sin h_0 = \sin\varphi\sin\delta - \cos\varphi\cos\delta ,
\]
que la règle applique sans la dire. Sur l'horizon droit ($\varphi = 0$),
$\cos H = \sin h\sec\delta$, et si de plus le Soleil est à l'équateur,
$H = 90^\circ - h$. Il rappelle enfin la règle d'arrondi : ajouter une unité si
les figures rejetées valent plus de $50\,000$.

Le volume se termine, juste après, sur la ligne du copiste : « Compendium
Mathematicæ. omittitur quia a P. M. Mersenno editum in Apologia scientiarum
cap. ultimo. »\footnote{« Le \emph{Compendium Mathematicæ} est omis, parce que
le P. M. Mersenne l'a publié dans l'\emph{Apologia scientiarum}, au dernier
chapitre. » La ligne « Compendium Mathematicæ » est centrée comme un titre ;
« omittitur quia … » est ajouté à sa suite, de la même main et de la même encre ;
le reste de la page est blanc, et la page de droite de la vue 49 aussi. Le cachet
de la Bibliothèque impériale est dans la marge (vue 49, gauche).}

\section*{Le premier livre des Sphériques de Maurolico : la fin}

\pagerange{20}{31}
À partir de la page de droite de la vue 20, et sans interruption jusqu'à la
vue 47, le volume donne la trigonométrie sphérique de Maurolico lui-même. Le
début du premier livre manque : le texte s'ouvre sur la fin d'une démonstration,
et la première proposition entière est la XXIII.\footnote{Le titre du premier
livre n'est pas dans le volume ; le second livre le cite comme « præcedentis
libri » ou « præmissi libelli ». La transcription du lot 1 note que cette
numérotation n'est pas celle de Théodose. Au premier livre appartiennent
encore les propositions 6, 7, 13, 16, 17 et 22 que citent les démonstrations,
et qui ne sont pas dans le volume. Cachet de la Bibliothèque impériale en marge
de la proposition XXIIII (vue 20, droite, numéro à l'encre 19).}

\emph{Notations.} Un triangle sphérique $ABC$ est formé d'arcs de grands cercles
plus petits qu'un quadrant ; on note $a = BC$, $b = CA$, $c = AB$ les côtés, et
$A$, $B$, $C$ les angles. Quand le triangle est rectangle, l'angle droit est en
$C$ et $c$ est l'hypoténuse.\footnote{Le feuillet dit « triangulum ex arcubus
circulorum majorum in superficie sphæræ » : « triangle d'arcs de grands cercles
sur la surface de la sphère », notre triangle sphérique ; il nomme les points par
des minuscules, qu'on met ici en capitales, et surmonte souvent les lettres d'un
trait, qu'on ne reproduit pas.} Maurolico raisonne sur une figure fixe, qu'il
appelle « la description de la proposition 17 » : on prolonge $AC$ et $AB$
jusqu'aux points $D$ et $E$ tels que $AD = AE = 90^\circ$, on mène le grand cercle
$DE$, et l'on prolonge $CB$ jusqu'à sa rencontre $F$ avec lui. Alors $F$ est le pôle
du cercle $AC$, de sorte que $FC = FD = 90^\circ$ ; l'arc $DE$ mesure l'angle
$A$ ; et
\[
FE = 90^\circ - A,\qquad FB = 90^\circ - a,\qquad BE = 90^\circ - c,\qquad
CD = 90^\circ - b .
\]
Chaque énoncé en « sinus secondaires » se lit donc en cosinus : le sinus de $FB$
est $\cos a$, celui de $FE$ est $\cos A$, et ainsi de suite. Les relations que
les démonstrations tirent des propositions 6 et 7 du livre sont, en termes
modernes, $\cos c = \cos a\cos b$, $\cos A = \cos a\sin B$ et
$\cos a/\cos A = \sin c/\sin b$, trois des relations du triangle rectangle
qu'on résume depuis le XVII\textsuperscript{e} siècle par les règles de
Napier.\footnote{Les règles de Napier sont postérieures au feuillet ; on les nomme
pour situer les énoncés, non pour dire que Maurolico les connaissait sous cette
forme. Les numéros 6 et 7 sont ceux des renvois ; le contenu qu'on leur prête est
celui que les démonstrations en tirent.}

\subsection*{La fin d'une proposition, et deux lemmes plans}

\pagerange{20}{20}
La page s'ouvre sur la conclusion d'une proposition dont l'énoncé manque : si
$BD$ est l'arc perpendiculaire mené de $B$ sur $AC$, alors
\[
\frac{\sin \widehat{ABD}}{\sin \widehat{CBD}} = \frac{\cos A}{\cos C},
\]
ce qui se tire des deux triangles rectangles $ABD$ et $CBD$, où
$\cos A = \cos BD \sin\widehat{ABD}$ et $\cos C = \cos BD \sin\widehat{CBD}$ ;
c'est bien l'étape qui reste écrite.\footnote{« sicut sinus totus ad sinum
anguli $cbd$ sic sinus secundus arcus $bd$ ad sinum secundum anguli $c$. Igitur ex
æquali sicut sinus anguli $abd$ ad sinum anguli $cbd$ sic sinus secundus anguli
$a$ ad sinum secundum anguli $c$. » La proposition XXVI du second livre (vue 42)
invoque « la 22 du livre précédent » pour exactement cet énoncé ; que la page
s'ouvre sur la fin de la proposition XXII est un rapprochement de cette lecture.}

Suivent deux propositions de géométrie plane, qui servent de lemmes.
\emph{XXIII} : deux triangles rectilignes de même aire et de mêmes angles sont
égaux (isométriques). \emph{XXIIII} : deux triangles de même aire construits sur une
même base, du même côté, et ayant même angle au sommet, ont leurs autres côtés
égaux deux à deux. C'est vrai : les sommets sont sur une même parallèle à la base
et sur un même arc capable, donc symétriques par rapport à la médiatrice de la
base ; la preuve du feuillet, qui renvoie à XXIII, suppose cette symétrie sans la
dire.\footnote{Le feuillet nomme les deux triangles « $abc$ $eda$ » et conclut
« latera $ab$ $ed$ erunt æqualia » ; la figure est un quadrilatère $d$, $b$ en
haut, $a$, $c$ en bas, dont les diagonales se coupent en $e$, et le second
triangle est $cda$, avec $AB = CD$ et $BC = DA$. « æqualia » est écrit deux fois.}

\subsection*{Propositions XXV à XXXII : le cas où les deux côtés de l'angle font un quadrant}

\pagerange{20}{23}
Fixons l'angle aigu $A$ et considérons le triangle rectangle $ABC$. Les huit
propositions portent sur quatre conditions et en établissent l'équivalence :

\begin{enumerate}
  \item[(i)] $\cos^2 a = \cos A$ ;
  \item[(ii)] $b + c = 90^\circ$ — ou, ce qui revient au même, $AB = CD$ et
    $AC = BE$ ;
  \item[(iii)] $\cos a = \sin B$ ;
  \item[(iv)] parmi tous les triangles rectangles de même angle $A$, la
    différence $c - b$ entre l'hypoténuse et le côté adjacent est la plus
    grande.
\end{enumerate}

XXV : (i) entraîne (ii) et (iii). XXVI : (i) entraîne (ii). XXVII : (ii)
entraîne (i) et (iii). XXVIII : (iii) entraîne (i) et (ii). XXIX : (i) entraîne
(iv). XXX : (iv) entraîne (i) et (ii). XXXI : (ii) entraîne (iv). XXXII : (iv)
entraîne (iii).\footnote{L'énoncé XXV est écrit « sinus unus secundus arcuum
acutis oppositorum fuerit medius proportionalis inter sinum totum sinumque acuti
anguli oppositi », « secundus » au-dessus de « unus » non biffé ; la démonstration
emploie le « sinus secondaire » de l'angle, et c'est la lecture retenue. Dans la
démonstration, l'hypothèse est écrite de façon incomplète, « Supponatur sinus
arcus $fb$ (qui est sinus secundus arcus $cd$ hoc est anguli $bac$ », sans
fermeture de la parenthèse ni second membre : $\sin FB$ est $\cos a$, et c'est
$\sin FE$ qui est le cosinus de l'angle $A$ (vue 21, gauche).}

Toutes ces implications sont vraies. En termes modernes, (i) équivaut à (iii)
par $\cos A = \cos a\sin B$ ; et avec $k = \cos A$ et $\tan b = k\tan c$,
\[
\tan(c-b) = \frac{(1-k)\tan c}{1 + k\tan^2 c},
\]
qui est maximal pour $\tan c = 1/\sqrt k$, c'est-à-dire $\tan b\tan c = 1$,
c'est-à-dire $b + c = 90^\circ$ ; alors $\cos a = \cos c/\cos b = \tan b$, d'où
$\cos^2 a = \tan^2 b = k$. La plus grande différence vaut
$\sin(c-b) = (1-k)/(1+k) = \tan^2(A/2)$ ; le premier livre ne la calcule pas, mais
elle suit des propositions XVI et XX du second.\footnote{Ce calcul est de cette
lecture. Pour la proposition XXX, qui raisonne par l'absurde (« si le sinus de
$FB$ n'est pas moyen proportionnel, soit $FG$ celui qui l'est »), il faut que le
triangle réalisant (i) existe : il existe toujours, $\tan b = \sqrt{\cos A}$
définissant un $b$ entre $0^\circ$ et $45^\circ$.}

La démonstration de XXV mérite d'être suivie, parce qu'elle est d'une autre
facture que la nôtre. De (i) et des propositions 6 et 7, Maurolico tire
$\sin c : \sin b = \cos b : \cos c$, soit $\sin 2c = \sin 2b$. Il double les arcs,
remplace les sinus par les demi-cordes, construit deux triangles rectilignes
inscrits dans des demi-cercles égaux, de même aire et de même angle droit au
sommet, et conclut par le lemme XXIIII que les cordes sont égales deux à deux,
« la plus grande à la plus grande ». Cela donne $2c = 180^\circ - 2b$ ; l'autre
possibilité, $c = b$, est exclue parce que l'hypoténuse est plus grande que le
côté ($\cos c = \cos a\cos b < \cos b$), ce que la clause « la plus grande à la
plus grande » contient.\footnote{« Quandoquidem chordæ duplatorum arcuum sunt
dupla sinuum simplis arcubus debitorum per diff. » ; « permutatam » et
« chordæ » (dans « chorda $gb$ chordæ $cl$ ») sont des lectures incertaines ;
« arcus » est ajouté dans l'interligne ; « chorda chorda » et « sinus sinus » sont
écrits deux fois ; le long d'un arc de la première figure est écrit « Inutilis »
(vues 21, gauche et droite). Aux numéros XXVI et XXVII, le dernier I est noyé
sous une tache : « XXVII » corrigé en « XXVI », et de même au suivant.}

La démonstration de XXIX compare de proche en proche les différences pour quatre
autres triangles de même angle, deux de chaque côté du triangle $ABC$, en menant
chaque fois un arc perpendiculaire au cercle $DE$. Elle est juste, mais un passage
y est copié deux fois en se brouillant.\footnote{De « Adhuc major erit » à « Quod
est secundum » (vue 22, droite), le copiste répète en le brouillant le passage
précédent ; à la vue 23, gauche, une phrase est écrite deux fois de suite ; dans
la dernière partie, un passage biffé (« $kp$ ad sinum arcus $qy$. Igitur major
sinus arcus $kp$ ») précède la bonne leçon. « dria » abrège « differentia »
(vue 22, gauche) ; « ut. inde » est une lecture incertaine. À XXX, un mot biffé
non lu, avec au-dessus un mot court lui aussi barré ; « medius » a la fin
surchargée ; à XXXI, « differrent » est incertain ; à XXVII, « Iter » est écrit
pour « inter ».}

\subsection*{Propositions XXXIII à XLVIII : deux points et leurs perpendiculaires}

\pagerange{23}{30}
Seconde série. Sur le côté $AE$ de l'angle aigu $A$, on marque deux points $B$ et
$G$, et l'on mène de chacun l'arc perpendiculaire au côté $AD$, de pieds $C$ et
$H$ ; ces arcs passent par $F$, et $FB = 90^\circ - BC$, $FG = 90^\circ - GH$.

\emph{XXXIII et XXXIV.}
\[
\frac{\sin CH}{\sin BG} = \frac{\cos A}{\cos BC\,\cos GH} .
\]
XXXIII l'écrit comme composition de deux rapports, $R : \sin FG$ et
$\sin FE : \sin FB$ ; XXXIV comme rapport de deux rectangles, celui du rayon et
du sinus de $FE$ à celui des sinus de $FG$ et $FB$. On l'a vérifié
numériquement.\footnote{À XXXIII, le feuillet écrit « sicut sinus arcus $gs$ ad
sinum arcus $bg$ : sic sinus arcus $fe$ ad sinum arcus $fg$ », là où la
conclusion porte, à juste titre, « ad sinum arcus $fb$ » (vue 24, gauche) ;
« signata » et « Hic » (ou « Sic ») sont des lectures incertaines ; les fins de
ligne de la vue 24 s'engagent dans la reliure. La démonstration invoque « 24. 6.
Element. », la composition des rapports des côtés dans les parallélogrammes
équiangles.}

\emph{XXXV.} Si $\cos A = \cos BC\,\cos GH$, alors $CH = BG$ ; $AG + AC = 90^\circ$
et $AH + AB = 90^\circ$ (« les portions alternes des quadrants sont égales »,
$AG = CD$ et $AH = BE$) ; enfin $\sin\widehat{ABC} = \cos GH$ et
$\sin\widehat{AGH} = \cos BC$. Tout cela est exact, et vérifié
numériquement. La démonstration de l'égalité $AG = CD$ passe, comme à XXV, par
des cordes d'arcs doubles et le lemme XXIIII.\footnote{Deux lettres fautives : « quadrangulum ex sinu toto
sinuque arcus $fg$ » où la proportion posée demande $fe$, et « sicut sinus arcus
$gf$ ad sinum arcus $be$ » où la suite demande $cd$ (vue 24, droite). « $km$ » est
écrit dans l'interligne (vue 25, gauche).}

\emph{XXXVI à XXXVIII} sont les réciproques : chacune des conditions $CH = BG$,
$AG = CD$ (ou $AH = BE$), $\sin\widehat{ABC} = \cos GH$ (ou
$\sin\widehat{AGH} = \cos BC$) entraîne $\cos A = \cos BC\cos GH$ et tout le
reste. XXXVII procède par l'absurde, en menant un petit cercle de pôle $F$.\footnote{À
XXXVII, « Ponatur $bg$ arcus æqualis arcui $cd$ » est écrit ainsi, pour
$ag$ ; un passage de trois lignes (« Quamobrem per 30 erit et arcus $ch$ … sinui
arcus $fg$ ») est biffé, puis réécrit plus bas avec le bon renvoi à 35 (vue 26,
gauche).}

\emph{XXXIX et XL.} Si $\cos A > \cos BC\,\cos KL$ (le point $K$ remplaçant $G$),
alors $CL > BK$, $AK > CD$, $AL > BE$, $\sin\widehat{AKL} > \cos BC$ et
$\sin\widehat{ABC} > \cos KL$ ; si $\cos A < \cos BC\,\cos KL$, toutes ces
inégalités sont renversées. La démonstration ramène chaque cas au cas d'égalité
en construisant le point $G$ de XXXV par un petit cercle de pôle $F$. Les deux
énoncés sont vrais ; on les a vérifiés sur un grand nombre de configurations
tirées au hasard.\footnote{L'énoncé de XL écrit « Et arcus ad angulum majores
erunt singuli singulis coalternis », alors que la démonstration conclut, à juste
titre, « arcus $ak$ minor arcu $cd$ » (vue 27, droite). Dans la démonstration de
XL, une phrase est écrite trois fois de suite ; « 16. 6. Eucl. » est une lecture
incertaine, où il faut « 15. 6. » comme à XXXIX ; « minor arcu $bd$ » est biffé
(vue 28, gauche). Le copiste passe de « XXXXVI » à « XLVII ».}

\emph{XLI à XLVI} sont les réciproques de XXXIX et XL, démontrées par trichotomie :
si $CL > BK$, ou $AK > CD$, ou $\sin\widehat{AKL} > \cos BC$, alors
$\cos A > \cos BC\cos KL$, puisque l'égalité contredirait XXXV et l'inégalité
contraire XL ; et de même dans l'autre sens.\footnote{À XLII, le cas d'égalité est
renvoyé « per 30 huius », où c'est la proposition 35 qui sert. À XLI, un mot biffé
non lu, avec « est » dans l'interligne ; à XLIV, « Q̄ est hypothesim » abrège
« quod est [contra] hypothesim », « contra » manquant ; à XLV, « minor » est ajouté
dans l'interligne (vues 28 et 29).}

\emph{XLVII.} Si deux couples de points $(G, L)$ et $(M, Q)$ sur $AE$ vérifient
$\cos GH\,\cos LK = \cos A$ et $\cos MN\,\cos QP = \cos GH\,\cos LK$ (où $H, K, N, P$
sont les pieds des perpendiculaires), alors $MG = KP$ et $NH = LQ$. En effet, par
XXXV appliquée aux deux couples, $AG = KD$, $AM = PD$, $AH = LE$, $AN = QE$, et
les différences donnent le résultat.\footnote{Le feuillet écrit « per 34 huius
arcus $am$ $pd$ æquales erunt », où c'est la proposition 35 qui donne l'égalité ;
la phrase suivante renvoie bien à 35 (vue 30, gauche). La permutation des rapports
est justifiée par « 23. 5. », les rapports « perturbés ».}

\emph{XLVIII.} Si de plus $\cos^2 BC = \cos GH\,\cos LK$, le point $B$ étant entre
les deux, alors $BG = CK$ et $HC = BL$ : $B$ vérifie la condition (i), donc
$AB = CD$ par XXV, et $AG = KD$ par XXXV.

\subsection*{Proposition XLIX : une autre démonstration de XXIX}

\pagerange{30}{31}
« Ce que la 29\textsuperscript{e} a proposé, le montrer autrement. » Si
$\cos^2 BC = \cos A$ et si $G$ est entre $A$ et $B$, alors $\cos GH > \cos BC$,
donc $\cos A = \cos^2 BC < \cos BC\cos GH$, et XXXIV donne $\sin HC < \sin BG$,
d'où $HC < BG$, c'est-à-dire $AB - AC > AG - AH$. De même pour un point au-delà
de $B$. La démonstration est plus courte que celle de XXIX et
juste.\footnote{« Minus autem quadratum sinus arcus $fb$ quadrangulo ex sinibus
arcuum $fb$ $bg$ producto » : il faut $fb$ $fg$, comme à la phrase suivante. « ad
quadrangulum » est écrit deux fois ; « qu » biffé ; la page s'arrête sur
« propon: », et la vue 31 reprend « ebatur demonstrandum » (vues 30, droite, et
31, gauche).}

\section*{Le second livre des Sphériques de Maurolico}

\pagerange{31}{40}
« Maurolyci Siculi Sphæricorum liber secundus » s'ouvre sur une préface qui dit
son sujet : le rapport du sinus de la somme au sinus de la différence des deux
arcs qui comprennent l'angle aigu d'un triangle sphérique rectangle. Ménélas,
dit Maurolico, ne l'a pas omis — c'est la cinquième proposition de son troisième
livre —, « mais nous, démontrant ce très noble théorème autrement et encore
autrement, l'avons éclairé de beaucoup de spéculations et comme de corollaires,
sans épargner les préambules utiles ».\footnote{« Qui locus quamvis a Menelao
minime sit prætermissus, Nos tamen Theorema illud nobilissimum quod ipsi quintum
est in ordine propositionum tertij libelli, aliter atque aliter demonstrantes
multis rem speculationibus Et quasi corollarijs illustravimus » (vue 31,
gauche). L'attribution à Ménélas est celle du feuillet ; cette lecture ne la
vérifie pas. « demonstratõem » abrège « demonstrationem ».} Le livre compte
trente-deux propositions ; les vingt-trois premières sont aux vues 31 à 40, la
suite aux vues 41--42 et 45--47.

\subsection*{I à V : le théorème de Ménélas}

Deux arcs $BA$ et $BC$ partent d'un même point $B$ ; deux autres, $CD$ et $AE$,
menés de leurs extrémités aux points $D$ de $BA$ et $E$ de $BC$, se coupent en $F$.
Les propositions I à V établissent, de deux manières chacune,
\[
\text{(I, II, III)}\quad
\frac{\sin AB}{\sin BD} = \frac{\sin AE}{\sin EF}\cdot\frac{\sin FC}{\sin CD},
\qquad
\text{(IV, V)}\quad
\frac{\sin BA}{\sin AD} = \frac{\sin BE}{\sin EC}\cdot\frac{\sin CF}{\sin FD},
\]
II n'étant qu'un réarrangement de I. La première égalité est le théorème de
Ménélas pour le triangle $ADF$ coupé par la transversale $BEC$, la seconde pour
le triangle $BDC$ coupé par la transversale $AFE$ : ce sont les deux formes du
théorème du quadrilatère complet sphérique, et elles sont justes (vérifiées
numériquement).\footnote{Le feuillet dit « rapport composé » là où l'on écrit un
produit de quotients. À I, l'énoncé de la démonstration porte « ex ratione sinus
arcus $fe$ ad sinum arcus $dc$ », la fin de la démonstration « $fc$ », qui est
juste (vue 31, droite). À IV, « duos ab eorum terminis ad eosdem reflexos $mf$
$ab$ » où les arcs de la figure sont $md$ et $ae$ ; à V, « arcus $bc$ » pour $bg$,
et « sicut arcus $cf$ » sans « sinus » (vue 32).} Les démonstrations de I, III et
V abaissent des arcs perpendiculaires sur un côté et emploient la proposition 7 du
premier livre — dans deux triangles rectangles de même angle, les sinus des
hypoténuses sont comme les sinus des côtés opposés à cet angle — ; celle de II
est un calcul de rapports avec une ligne auxiliaire $l$ ; celle de IV prolonge
$BC$ et $DC$ jusqu'à leur point de rencontre $M$, qui est à un demi-cercle, et
remplace chaque arc par son supplément, qui a même sinus.

\subsection*{VI à IX : les lemmes plans}

\pagerange{33}{34}
\emph{VI}, le théorème de Ptolémée : dans un quadrilatère inscrit $ABCD$,
$AC\cdot BD = AB\cdot CD + AD\cdot BC$, démontré par la construction classique d'un
point $E$ de la diagonale tel que $\widehat{DAE} = \widehat{BAC}$.

\emph{VII} : pour deux arcs inégaux $\alpha > \beta$ d'un même cercle, et en notant
$\operatorname{crd}$ la corde,
\[
\frac{\operatorname{crd}\alpha\,\operatorname{crd}(180^\circ-\beta) + \operatorname{crd}\beta\,\operatorname{crd}(180^\circ-\alpha)}
     {\operatorname{crd}\alpha\,\operatorname{crd}(180^\circ-\beta) - \operatorname{crd}\beta\,\operatorname{crd}(180^\circ-\alpha)}
= \frac{\operatorname{crd}(\alpha+\beta)}{\operatorname{crd}(\alpha-\beta)} ,
\]
par Ptolémée appliqué à deux quadrilatères.

\emph{VIII}, la même chose en sinus :
\[
\frac{\sin p\cos q + \cos p\sin q}{\sin p\cos q - \cos p\sin q} = \frac{\sin(p+q)}{\sin(p-q)} .
\]
La construction est belle : sur le rayon $BD$ comme diamètre on décrit un cercle,
qui passe par les pieds $A$, $C$ des perpendiculaires menées de $B$ sur les rayons
$DP$, $DQ$ et par le point $E$ symétrique ; dans ce petit cercle, $BA = \sin p$,
$BC = BE = \sin q$, $DA = \cos p$, $DC = DE = \cos q$, $AC = \sin(p+q)$ et
$AE = \sin(p-q)$, et VII s'applique. La figure donne en fait davantage que ce
que le feuillet énonce : le théorème de Ptolémée appliqué au quadrilatère inscrit
$BADC$ (resp. $BAED$) donne séparément
$\sin(p\pm q) = \sin p\cos q \pm \cos p\sin q$, les formules d'addition ; le
feuillet n'en tire que le rapport.\footnote{Remarque de cette lecture. Le feuillet
écrit « lineæ perpendiculares $ba$ $bc$ » là où la figure porte $e$ au pied de la
seconde, et où la suite dit « linea $bc$ vel linea $be$ » ; la seconde lettre de
« $ae$ » ressemble à un c ; un « erit » biffé (vue 34, gauche).}

\emph{IX} : si $a : b = c : d$ (avec $a > b$, $c > d$), alors
$(a+b) : (a-b) = (c+d) : (c-d)$, et réciproquement.\footnote{« Aio quod erit $ab$
ad $bc$ » : le mot « ad » est écrit en début de ligne comme « $ab$ » (vue 34,
droite).}

\subsection*{X à XIV : le théorème de la somme et de la différence, trois fois}

\pagerange{35}{37}
\emph{X.} Dans un triangle sphérique rectangle en $C$, dont $A$ est aigu,
\[
\frac{\sin(c+b)}{\sin(c-b)} = \frac{1+\cos A}{1-\cos A} .
\]
C'est exact : $\tan b = \tan c\cos A$ s'écrit $\sin b\cos c = \cos A\sin c\cos b$,
et il suffit de former somme et différence de $\sin c\cos b$ et
$\cos c\sin b$. C'est exactement le chemin de Maurolico : les deux rectangles
$g = \sin c\cos b$ et $h = \sin b\cos c$ sont dans le rapport $1 : \cos A$ par la
proposition IV, IX donne $(g+h):(g-h) = (1+\cos A):(1-\cos A)$, et VIII identifie
$(g+h):(g-h)$ à $\sin(c+b):\sin(c-b)$.\footnote{Dans la « description de la
dixième », les quadrants $AD$ et $AE$ sont pris sur $AB$ et $AC$ prolongés, et $F$
est la rencontre de $DE$ et de $CB$ ; le sinus de $DE$ est celui de l'angle $A$, le
sinus de $FE$ son cosinus. La figure est celle du premier livre, les lettres $D$ et
$E$ échangées.}

\emph{XI} ramène le rapport à une figure plane : si l'on construit en $H$ deux
angles $\widehat{GHK} = c$ et $\widehat{GHL} = b$ de part et d'autre d'une droite
$HG$ perpendiculaire à $LK$, alors $1 : \cos A = KG : GL$ — ce qui est
$\tan c : \tan b$. \emph{XII} en tire une seconde démonstration de X : avec
$HK$ égal au rayon et $GM = GL$, les perpendiculaires $KN$ et $KO$ sont
$\sin(b+c)$ et $\sin(c-b)$, et des triangles semblables donnent
$KN : KO = LK : KM = (KG+GL) : (KG-GL)$. Une scolie distingue trois figures,
selon que $b + c$ est plus petit, plus grand qu'un quadrant, ou égal, et note que
la même démonstration sert aux trois ; elle annonce une troisième méthode, « que
j'estime transmise par Ménélas ».\footnote{« superest nunc tertius idipsum
demonstrandi modus quem a Menelao traditum existimo » (vue 36, droite). Dans
XII, la démonstration que $c > b$ renvoie à « 28. 2. sphæric. elementorum ».}

\emph{XIII}, un lemme : si un rayon coupe en $D$ une corde $AC$ et en $B$ son arc,
$AD : DC = \sin AB : \sin BC$. \emph{XIV}, la troisième démonstration, est dans
l'espace : le petit cercle de pôle $A$ passant par $B$ coupe le grand cercle $AC$
en $D$ et $E$, de sorte que $CE = b + c$ et $CD = c - b$ ; la perpendiculaire
$BG$ au plan du grand cercle $AC$ tombe sur la corde $DE$ et sur le rayon $HC$
de la sphère ; dans le petit cercle de centre $F$, l'angle $\widehat{BFD}$ est
l'angle $A$ du triangle, de sorte que, le rayon $BF$ étant pris pour unité,
$EG = 1 + \cos A$ et $GD = 1 - \cos A$ ; et XIII appliquée au rayon $HC$ donne
$\sin CE : \sin CD = EG : GD$. C'est une démonstration complète.\footnote{« Itaque
demonstrandum est quod arcus $cae$ ad sinum arcus $cd$ » : « sinus » manque
devant « arcus $cae$ » ; « \emph{secundum} » est biffé dans « linea $hgc$
semidiameter », et « ut » dans « Et hoc ut antea » (vue 37). La perpendicularité
de $BG$ est tirée de « 19. 11. Euclidis », celle de $AF$ au plan du petit cercle
de « 10. 1. Elem. Theodosij ».}

\subsection*{XV à XIX : demi-angle et sinus verses}

\pagerange{37}{39}
\emph{XV} : $(1+\cos A) : (1-\cos A) = \cos^2(A/2) : \sin^2(A/2)$, par un
demi-cercle de diamètre $HK$ où l'arc $KM$ vaut $A$ et $MN$ est la perpendiculaire.
\emph{XVI} : donc
\[
\frac{\sin(c+b)}{\sin(c-b)} = \cot^2\frac{A}{2} .
\]
\emph{XVII} : c'est aussi le rapport $\operatorname{vers}(180^\circ - A) :
\operatorname{vers} A$, le « sinus verse majeur » de l'angle à son sinus verse.
Les trois énoncés sont exacts.\footnote{À XV, la démonstration écrit « sicut
sinus arcus $fg$ ad quadratum sinus arcus $gd$ », puis « sicut quadratum sinus
arcus $fg$ ad sinum arcus $gd$ », chaque fois avec un « quadratum » mal placé
(vue 38, gauche). À XVI, « arcuum $bac$ $ac$ », la lettre du milieu surchargée. À
XVII, « sicut linea $nk$ ad lineam $nk$ », le k repassé d'une encre plus noire, où
le sens demande $hn$ ; un « $hk$ » biffé et douteux.}

\emph{XVIII} : deux triangles rectangles qui ont un angle aigu égal ont même
rapport $\sin(c+b) : \sin(c-b)$. \emph{XIX} : réciproquement, si ces rapports sont
égaux, les angles aigus sont égaux — puisque, la somme des deux sinus verses
étant le diamètre, l'égalité des rapports $\operatorname{vers}(180^\circ-A) :
\operatorname{vers}A$ et $\operatorname{vers}(180^\circ-D) : \operatorname{vers}D$
donne, par composition, $\operatorname{vers}A = \operatorname{vers}D$. Vrai ; la
copie de XIX est pourtant incomplète dans l'énoncé et brouillée dans la
démonstration.\footnote{L'énoncé passe, au changement de page, de « ad sinum
arcus differentiæ eorumdem » à « tum ipsi acuti anguli æquales erunt », sans le
second membre de la proportion. Dans la démonstration, la comparaison est réduite
à un seul membre (« sicut aggregatum sinuum versorum anguli acuti $d$ ad sinum
versum anguli $d$ »), celui qui concerne l'angle $A$ manquant ; « etiam idem » est
biffé, avec « tam illud » dans l'interligne, et « diameter » est suppléé par la
transcription (vues 38 et 39). À XVIII, deux passages biffés, « $ed$ $df$ » et
« sinibus versis anguli ».}

\subsection*{XX et XXI : la plus grande différence, deux fois encore}

\pagerange{39}{40}
\emph{XX.} Dans un triangle rectangle, l'angle aigu $A$ étant fixé, la différence
$c - b$ est la plus grande quand $b + c = 90^\circ$, et seulement alors. La
démonstration est la plus directe de toutes : par XVIII,
$\sin(c-b) = \sin(c+b)\,\tan^2(A/2) \le \tan^2(A/2)$, avec égalité si et
seulement si $\sin(c+b) = 1$ ; et $c - b$ étant compris entre $0^\circ$ et
$90^\circ$, maximiser son sinus revient à maximiser l'arc. On retrouve les
propositions XXIX à XXXII du premier livre, avec la valeur du maximum,
$\sin(c-b) = \tan^2(A/2)$, que le feuillet ne tire pas.\footnote{La réciproque est
donnée deux fois, directement et par l'absurde ; une phrase y est écrite deux fois
de suite ; « propositis », en fin de démonstration, est abrégé et incertain. Sur
la figure, l'initiale « $k$ » ressemble à un x (vues 39, gauche et droite).}

\emph{XXI.} Autre démonstration, dans la figure plane de XII, où
$\widehat{KHO} = c - b$ et $\widehat{LHK} = b + c$ : l'angle $\widehat{KHO}$ est
maximal quand le cercle passant par $H$, $M$ et $K$ est tangent en $H$ à la droite
$GH$, c'est-à-dire quand $GH^2 = GM\cdot GK = GL\cdot GK$, c'est-à-dire quand
$\widehat{LHK}$ est droit. C'est le critère classique de l'angle maximal sous
lequel on voit un segment depuis une droite, et la démonstration est
juste.\footnote{« per 26. 3. Et 26. 1. Eucli. » : les deux nombres sont tracés de
même, et le second est lu 26 sans certitude (vue 40, gauche).}

\subsection*{XXII à XXV : la règle des sinus pour deux triangles}

\pagerange{40}{41}
\emph{XXII.} Si deux triangles sphériques $ABC$ et $DEF$ ont $A = D$ aigus, et
$C + F = 180^\circ$ (ou $C = F = 90^\circ$), alors
$\sin AB : \sin DE = \sin BC : \sin EF$. En termes modernes, c'est la règle des
sinus, $\sin a/\sin A = \sin c/\sin C$, appliquée aux deux triangles, avec
$\sin C = \sin F$ ; la démonstration de Maurolico rabat l'angle obtus par
symétrie pour se ramener à deux angles égaux.\footnote{« per 7. \emph{huius}
præcedentis libri » : « huius » biffé (vue 40, droite).}

\emph{XXIII}, la réciproque : si $A = D$ aigus et $\sin AB : \sin DE = \sin BC :
\sin EF$, alors $C = F$ ou $C + F = 180^\circ$. C'est vrai, puisque la règle des
sinus donne $\sin C = \sin F$. La démonstration traite le cas où $C$ est droit,
puis mène les perpendiculaires $BH$ et $EG$ et distingue deux cas : toutes deux
tombent à l'intérieur des triangles, ou l'une dedans et l'autre dehors. Il manque
le cas où toutes deux tombent au dehors, c'est-à-dire où $C$ et $F$ sont tous deux
obtus ; il se traite exactement comme le premier et donne $C = F$.\footnote{L'angle
$A$ étant aigu, le pied de la perpendiculaire issue de $B$ tombe sur le côté $AC$
si et seulement si $C$ est aigu. Remarque de cette lecture. La démonstration
commence à la vue 40, droite, et s'arrête sur « quod » ; la vue 41, gauche,
reprend sur « impossibile ». Sur la seconde figure, la lettre $f$ est écrite
deux fois.}

\pagerange{41}{42}
\emph{XXIIII}, la bissectrice : si l'arc $BD$ coupe en deux l'angle $B$ du
triangle $ABC$, alors $\sin AB : \sin BC = \sin AD : \sin DC$ — par XXII, les
angles en $D$ étant supplémentaires. \emph{XXV}, la réciproque, par XXIII : les
deux angles $\widehat{ABD}$ et $\widehat{DBC}$ sont égaux ou supplémentaires, et
ils ne peuvent être supplémentaires puisque leur somme est l'angle $B$ du
triangle.\footnote{À XXV, « ponatur arcus $ab$ ad sinum $bc$ » : « sinus »
manque devant « arcus » ; « Producatque » est une lecture incertaine ; « ad ad »
(vue 41).}

\emph{XXVI.} Si deux triangles $ABC$, $DEF$ ont $A = D$ et $C = F$, les
perpendiculaires $BG$ et $EH$ partagent les bases dans le même rapport des sinus :
$\sin AG : \sin GC = \sin DH : \sin HF$. Car
$\sin AG : \sin GC = (\sin AB : \sin BC)(\sin\widehat{ABG} : \sin\widehat{GBC})$,
le premier rapport vaut $\sin C : \sin A$ et le second $\cos A : \cos C$ (la
proposition de la vue 20). Vrai, les portions $AG$ et $GC$ étant comptées sur le
grand cercle $AC$ même si le pied tombe au dehors.\footnote{« per 16 præmissi » :
le 6 a la forme en « b » de cette main (vue 42, gauche). Les renvois sont à 13,
16 et 22 du premier livre.}

\emph{XXVII.} Dans un triangle rectangle en $C$,
$\sin B : 1 = \cos A : \cos a$. \emph{XXVIII.}
$\cos A : 1 = \cos a\,\sin B : 1^2$. Ce sont deux écritures de
$\cos A = \cos a\sin B$, l'une des règles de Napier.\footnote{À XXVII, la
démonstration conclut « sinus arcus $cf$ ad sinum arcus $fb$ sic sinus anguli
$abc$ ad sinum totum », où il faut $ef$ ($\sin FC$ est le rayon) ; la transcription
note que la main forme $c$ et $e$ presque de même et donne les lettres telles
qu'elles se lisent. Dans la marge de XXVIII, un cercle lettré avec la légende
« Refertur ad propositionẽ Sequentem ».}

\subsection*{XXIX à XXXII : la loi des cosinus, et les ombres}

\pagerange{45}{47}
La démonstration de XXVIII, interrompue au bas de la vue 42, s'achève à la vue
45.\footnote{La page de gauche de la vue 42 s'arrête sur « quadrangulum quod » ;
la page de droite de la vue 45 (numéro à l'encre 44) reprend en les répétant les
deux dernières lignes de la vue 42 et achève la proposition. Les trois feuillets
intermédiaires portent le texte sur Regiomontanus, lu dans la section suivante.
La reconstitution est celle de la transcription.}

\emph{XXIX.} Dans \emph{tout} triangle sphérique,
\[
\frac{\operatorname{vers} A}{\operatorname{vers} a - \operatorname{vers}(b-c)}
= \frac{1}{\sin b\,\sin c},
\qquad\text{soit}\qquad
\operatorname{vers} a = \operatorname{vers}(b-c) + \sin b\,\sin c\,\operatorname{vers} A .
\]
C'est la loi des cosinus sphérique, $\cos a = \cos b\cos c + \sin b\sin c\cos A$,
écrite en sinus verses ; on l'a vérifiée numériquement. La démonstration est dans
l'espace, et en voici l'idée dans nos termes. Le petit cercle de pôle $A$ passant
par $C$ a pour rayon $\sin b$ ; l'arc qu'il intercepte entre le grand cercle $AB$
et $C$ est semblable à l'arc du grand cercle polaire de $A$ qui mesure l'angle
$A$, de sorte que sa flèche vaut $\sin b\,\operatorname{vers} A$. Projetée sur
l'axe du point $B$, cette flèche est multipliée par $\sin c$ ; et ce qu'elle
mesure sur cet axe est la distance entre les projections de $C$ et du point de
$AB$ situé à la distance $b$ de $A$, c'est-à-dire
$\operatorname{vers} a - \operatorname{vers}(b-c)$.\footnote{Les lettres de la
démonstration sont souvent douteuses : un signe en forme de 2 après « majores
quidem » ; le point « $r$ » ; la perpendiculaire « ad $os$ » ; une lettre en forme
de c retourné, que la transcription note $\supset$ ; « $aqz$ », repassé, qui ne
nomme aucune ligne de la figure et où l'énoncé demande $ry$ ; « vel $qz$ » (vues
45, droite, et 46, gauche). On n'a donc pas suivi la démonstration lettre à
lettre ; l'idée qu'on en donne est celle que ses étapes lisibles portent
(« arcus $dl$ $ks$ sunt similes », « in similibus arcubus sinus sunt semidiametris
proportionales », triangles équiangles $kts$ et $aux$).}

\emph{XXX.} Pour un arc $x$ et un gnomon égal au rayon :
$\sin x : \cos x = 1 : \text{umbra recta}$, et
$\cos x : \sin x = 1 : \text{umbra versa}$ ; autrement dit, l'« ombre droite » est
$\cot x$ et l'« ombre verse » $\tan x$. \emph{Corollaires} : l'ombre droite d'un
arc est l'ombre verse de son complément ; à $45^\circ$ les deux ombres égalent le
gnomon ; le gnomon est moyen proportionnel entre les ombres de deux arcs
complémentaires, $\tan x\,\tan(90^\circ - x) = 1$.\footnote{« recta $ca$ sinus
rectus arcus $bc$ » : lecture incertaine, et c'est la perpendiculaire $ce$ qui
est le sinus (vue 46, gauche).}

\emph{XXXI.} Dans un triangle rectangle en $C$,
$1 : \sin b = 1 : \tan(90^\circ - A)\,\tan a$, soit
\[
\sin b = \tan a\,\cot A .
\]
\emph{Corollaire} : deux triangles rectangles qui ont un côté $b$ égal ont même
produit $\tan a\cot A$, et réciproquement. \emph{XXXII} :
$1 : \sin b = \tan A : \tan a$, la même relation, par le dernier corollaire de
XXX ; le feuillet ajoute qu'on en conclut tout pareillement pour les ombres
droites des compléments. Ce sont encore des règles de Napier, et elles sont
exactes.\footnote{À XXXI, « respondet » est offert pour un mot abrégé lu « rndet »
avec un trait ; la phrase « Igitur ex æqua proportione erit sinus arcus $da$ …
ad sinum arcus $ac$ » est écrite deux fois, la seconde avec « hoc est sinus
totus » en interligne. À XXXII, un mot abrégé d'environ cinq lettres après
« descriptione » n'est pas lu (vues 46 et 47). Dans la marge de la vue 47,
gauche, deux « $f$ » isolés le long de la reliure.}

\subsection*{La scolie : des démonstrations au calcul}

\pagerange{47}{47}
« Par ces démonstrations on peut résoudre toutes les questions qu'on a coutume
de poser sur les triangles sphériques, et qui regardent en astronomie les arcs
des cercles conçus dans le premier ciel » ; mais qui les possède doit savoir
ramener la théorie à la pratique et la démonstration au calcul, et Maurolico
donne des exemples.\footnote{« Sed nemo harum speculationum scientiam perfectam
habens nesciet theoriam ad praxim atque demonstrationem ad calculum reducere »
(vue 47, gauche).} Dans la figure de XXXI, que $ACD$ soit l'équateur, $ABE$
l'écliptique, $FED$ le colure, $A$ l'équinoxe, $DE = \varepsilon$, $B$ le lieu de
l'astre et $FBC$ son cercle de déclinaison. Alors $1 : \sin\varepsilon = \sin
\lambda : \sin\delta$, et le quatrième terme se calcule en multipliant le
deuxième par le troisième et en divisant par le premier : c'est la règle du
problème 1.\footnote{« sic \emph{sit} sinus arcus $ab$ » : « sit » incertain ;
« per 6. \emph{huius} præmissi libri » : « huius » biffé.} Puis, « si l'on prend
l'arc $FBC$ pour horizon droit », $AC$ est l'ascension droite de l'arc $AB$, et la
proposition 17 du premier livre, $\cos\lambda = \cos\delta\cos\alpha$, la
donne.\footnote{La phrase s'arrête au bas de la vue 47, gauche, sur « Si autem
sumatur arcus $fbc$ pro horizonte recto » et se poursuit en tête de la vue 42,
droite (numéro à l'encre 41) : « erit arcus $ac$ ascensio recta, respondens arcui
Zodiaci $ab$ ». Le feuillet écrit « sicut sinus totus ad sinum arcus $ac$ », où il
faut le sinus secondaire de $AC$, comme le montre la reformulation qui suit
(« sicut sinus totus ad sinum arcus $cd$ »).}

\section*{Les règles de Regiomontanus, et la table des sécantes}

\pagerange{42}{45}
Sans titre, à la suite de la scolie, un texte suivi explique « d'où dépendent »
les règles de calcul des tables générales de « Iohannes de Monteregio »,
Regiomontanus, « car cela n'est pas clair pour tous les esprits », puis démontre
la table bienfaisante et en donne l'usage.\footnote{Ce texte occupe les pages
numérotées à l'encre 41 à 43 (vue 42, droite, à vue 45, gauche) et s'achève en
tête de la vue 47, droite : la vue 45, gauche, s'arrête sur « vel quot zyphræ
sunt in numero », et la vue 47, droite, s'ouvre sur « sinus totius, Et habeo in
quotiente sinum arcus $cd$ quæsiti ». Reconstitution de la transcription. Un signe
en croissant, « ( », marque de paragraphe de cette main, ouvre plusieurs
alinéas (vues 43, 44, 45 et 47) ; il n'est pas un mot.}

\subsection*{La table des déclinaisons}

Soient l'écliptique $AB$, l'équateur $AC$, $B$ le point de l'écliptique qui a la
longitude de l'astre, $G$ l'astre sur le cercle de latitude $BG$ prolongé jusqu'à
l'équateur en $C$, et $GH$ sa déclinaison. La table donne d'abord l'arc $BC$
(notre \emph{Primus}), d'où, avec la latitude, l'arc $GC$ ; puis un « nombre
multiplicande » qui est le sinus de l'angle $\widehat{GCH}$ pour un rayon de
$100\,000$ (notre \emph{Secundus}) — égal au rayon au point solsticial, où cet
angle est droit. Dans le triangle rectangle $GHC$, $1 : \sin\widehat{GCH} = \sin
CG : \sin GH$ : on connaît trois termes, et le quatrième s'obtient en multipliant
le deuxième par le troisième et en divisant par le premier. C'est la règle de
Regiomontanus, et c'est notre problème 4.\footnote{« Recte igitur præcipit
Iohannes in problemate dum iubet multiplicari sinum arcus $gc$, in numerum
multiplicandum in tabula repertum » (vue 42, droite). Après « arcus », un mot
surchargé et empâté, commençant peut-être par « vi », n'est pas lu ;
« \emph{primum} » est biffé.}

Diviser par $100\,000$, c'est rejeter cinq figures ; et l'on ajoute une unité au
reste si les figures rejetées valent plus de $50\,000$, parce que la fraction
négligée dépasse alors la moitié. C'est l'arrondi au plus proche ; le cas exact
de $50\,000$ n'est pas discuté. Il importe peu, ajoute-t-il, que les sinus de
$CG$ et $GH$ soient pris pour un autre rayon que celui de l'angle : dans une
proportion, seuls le premier et le deuxième termes, et le troisième et le
quatrième, doivent être de même espèce.\footnote{« quandoquidem in quatuor
magnitudinibus proportionalibus, non necesse est omnes quatuor esse eiusdem
colligantiæ, sed primam cum secundâ et tertiam cum quartâ » (vue 43,
gauche). « \emph{negligitur} » est biffé et remplacé.}

\subsection*{La table des ascensions droites}

Dans la figure de la proposition XXXI, soient l'équateur $ACD$, le cercle de
latitude de l'astre $ABE$, son cercle de déclinaison $FBC$, et $FED$ le cercle
passant par les pôles de ces deux-là ; $B$ est l'astre, $BC$ sa déclinaison, $AC$
l'arc $\Delta$ du problème 5. La table donne la « racine des ascensions », arc
d'équateur du début du Bélier au cercle de latitude (notre \emph{Tertius}), et un
nombre multiplicande qui est « l'ombre verse de l'arc $FE$ multipliée par six »,
pour un gnomon de $100\,000$. Or la proposition XXXI donne
\[
\sin AC = \tan FE\,\tan BC ,
\]
avec $FE = 90^\circ - \theta$. Regiomontanus prend un rayon de $60\,000$ et un
gnomon de $100\,000$, dont le carré vaut $10^{10}$ : le sinus de $AC$ pour ce rayon
est $60\,000\cdot\tan FE\cdot\tan BC/10^{10} = (6\tan FE)\cdot\tan BC/10^{6}$, les
tangentes étant en parties du gnomon. On multiplie donc le nombre multiplicande
par la tangente de la déclinaison et l'on rejette six figures ; le canon dit
cinq, parce qu'une figure a déjà été ôtée au nombre multiplicande « pour alléger
le travail ». L'ascension droite est la racine augmentée ou diminuée de $AC$.
Le calcul est juste.\footnote{Le carré du gnomon est écrit en chiffres serrés,
lu « 10000000.000 » en comptant dix zéros, compte que la transcription tient
pour douteux ; plus bas, « 10000.000.000 », dix zéros encore, le deuxième chiffre
ayant la forme en « a » du zéro de cette main. « sæpius » est surchargé ;
« haben » est écrit pour « haberi » (vues 43, gauche et droite).}

Maurolico remarque que $AC$ s'obtient aussi sans nombre multiplicande ni table
des ombres : dans le triangle rectangle $GHC$ de la figure précédente, la
proposition 17 du premier livre donne $\cos GH : \cos GC = 1 : \cos CH$, soit
$\cos\Delta = \cos s/\cos\delta$ (notre seconde règle du problème 5). Si
Regiomontanus, « homme très perspicace », a ajouté le nombre multiplicande et la
table féconde, c'est pour éviter cette division, pénible quand le diviseur a des
chiffres quelconques, et tout faire par multiplication.

\subsection*{La différence ascensionnelle}

Dans la même figure, que $ACD$ soit l'équateur, $FBC$ l'horizon droit, $ABE$
l'horizon oblique et $FED$ le méridien : $FE$ est la latitude du lieu, $BC$ la
déclinaison de l'astre, $CA$ sa différence ascensionnelle, et XXXI donne
$\sin\Delta = \tan\varphi\,\tan\delta$. La règle de Regiomontanus — multiplier les
ombres verses de la latitude et de la déclinaison, puis par six, et rejeter six
figures — est la même opération pour un rayon de $60\,000$. « Ainsi est démontrée
la pratique de ces tables, et en même temps déclaré comment elles ont été
construites. »\footnote{« ideoque » est une lecture incertaine (vue 43, droite).}

\subsection*{La démonstration de la table bienfaisante}

\pagerange{44}{44}
« À l'imitation de la table féconde de Iohannes de Monteregio, nous avons fait
une autre table, que nous avons appelée bienfaisante, parce que par son bienfait
on procure de la facilité à certains calculs. » Soit un triangle rectiligne
$ABC$ rectangle en $B$, avec la hauteur $BD$ sur $AC$, et $AB$ égal au rayon et au
gnomon, $100\,000$ parties. Pour l'angle $A$, la table féconde donne $BC$
(l'ombre verse, $\tan A$), la table bienfaisante $CA$ (le « rayon » qui joint le
sommet du gnomon à l'extrémité de l'ombre, $\sec A$), la table des sinus $BD$
($\sin A$), et $DA$ est $\cos A$. Les triangles $ABD$, $BCD$ et $ACB$ étant
semblables,
\[
DA \cdot CA = AB^2,\qquad BC = \frac{AC\cdot BD}{AB},\qquad CA = \frac{AB\cdot BC}{BD},
\]
soit $\sec A = 1/\cos A$, $\tan A = \sec A\sin A$, $\sec A = \tan A/\sin A$ ;
diviser par $AB$ ou multiplier par $AB$, c'est rejeter ou ajouter cinq
zéros.\footnote{Plusieurs mots de la page sont soulignés sans raison apparente ;
la transcription ne reproduit pas le soulignement (vue 44, gauche).}

Quatre valeurs particulières. À $30^\circ$, $\tan A$ est la moitié de $\sec A$, et
$AC^2 = \tfrac43 AB^2$ : si $AB$ est le rayon d'une sphère, $AC$ est le côté du cube
inscrit dans cette sphère, « mais cela ne fait rien à l'affaire ». À $45^\circ$,
$AD = DB = DC$, l'ombre égale le gnomon, et $AC$ est le côté du carré inscrit dans
le cercle de rayon $AB$. À $60^\circ$, $AC = 2AB = 200\,000$. À $90^\circ$, l'ombre et
le rayon « s'en vont à l'infini », les deux droites étant parallèles. Tout cela
est exact, et les deux tables se tirent de la table des sinus.\footnote{Le
feuillet dit « latus cubi circumscripti a sphæra », « le côté du cube circonscrit
par la sphère », c'est-à-dire inscrit en elle : son côté est
$2R/\sqrt3 = R\sqrt{4/3}$, ce qui est bien $AC$. « lineam iam $ac$ » est écrit
ainsi (vue 44, gauche et droite).}

\subsection*{L'usage de la table bienfaisante}

\pagerange{44}{47}
Soient deux quadrants de grands cercles $ABE$ et $ACD$ inclinés l'un sur l'autre,
$F$ le pôle de $ACD$, et les quadrants $FED$ et $FBC$, de sorte que les angles en
$C$, $D$ et $E$ sont droits. Dans le triangle $ABC$, rectangle en $C$, la sécante
permet, dit Maurolico, de résoudre par une multiplication et une suppression de
cinq figures les quatre problèmes suivants.\footnote{« et a puncto $f$ polo scilicet
ipsius $ae$ » : $F$ est dit pôle de « $ae$ » ; les angles droits en $C$ et $D$ et
l'usage qui suit en font le pôle du cercle $ACD$. « apud $c$ $d$ $e$ puncta » est
écrit « apud cde puncta » ; « quoad » est écrit « qual », corrigé par « oa » en
interligne (vue 44, droite).}

\begin{enumerate}
  \item[1.] \emph{Connaissant $AB$ et $BC$, trouver $AC$} :
    $\cos AC = \cos AB\,\sec BC$. Le feuillet l'écrit comme proportion continue
    ($\sin FB$, rayon et sécante de $BC$) combinée à la relation « démontrée par
    Ménélas » $\sin FB : 1 = \sin BE : \sin CD$, c'est-à-dire
    $\cos BC = \cos AB/\cos AC$. Applications : l'ascension droite par la
    longitude et la déclinaison (problème 3) ; l'arc $\Delta$ du problème 5 par
    l'arc du cercle de latitude et la déclinaison ; la différence ascensionnelle
    par l'amplitude et la déclinaison (problème 7).\footnote{« figuras » est
    écrit au-dessus de « zyphras » biffé (vue 44, droite).}
  \item[2.] \emph{Connaissant l'angle $A$ (l'arc $DE$) et $BC$, trouver $AB$} :
    $\sin AB = \sin BC\,\sec FE = \sin BC/\sin A$. Applications : la longitude
    par la déclinaison et l'obliquité ; l'amplitude ortive par la déclinaison et
    l'angle de l'horizon avec l'équateur, complément de la latitude
    (problème 6).\footnote{« diff\textsuperscript{ia} » et
    « quæstionib\textsuperscript{9} » abrègent « differentia » et
    « quæstionibus ». La phrase « si ex angulo $a$ complemento latitudinis
    regionis et ex arcu $bc$ declinationis stellæ » est écrite deux fois de
    suite, sans biffure (vue 45, gauche).}
  \item[3.] \emph{Connaissant $AC$ et $CB$, trouver $AB$} : par la règle des
    déclinaisons appliquée aux compléments, $\cos AB = \cos AC\,\cos BC$.
  \item[4.] \emph{Connaissant $AB$ et $BC$, trouver l'angle $A$} :
    $\sin A = \sin BC\,\sec BE = \sin BC/\sin AB$ — ainsi « de l'arc $AB$ de
    l'écliptique et de sa déclinaison $BC$ nous obtiendrons la plus grande
    déclinaison $DE$ ».\footnote{Le feuillet appelle l'arc cherché « $cd$ » à deux
    reprises (« quæratur arcus $cd$ », « habeo in quotiente sinum arcus $cd$
    quæsiti »), puis conclut sur « maximam declinationem $de$ ». C'est $DE$ :
    la proportion qu'il écrit, $\sin AB : 1 = \sin BC : \sin(\cdot)$, est
    $\sin\delta = \sin\lambda\sin\varepsilon$. « prius » est surchargé, peut-être
    sur « primus » (vue 45, gauche, et vue 47, droite).}
\end{enumerate}

\subsection*{La date de 1530}

\pagerange{47}{47}
Le texte conclut que tout ce que Regiomontanus calculait par la table féconde,
évitant l'ennui de la division, se calcule aussi aisément par « notre » table
bienfaisante, et promet des exemples pratiques « plus bas, avec les tables » —
dans ce volume, où l'ordre du livre n'est pas celui de la reliure, les problèmes
résolus avec la table bienfaisante y répondent sans doute, sans que rien
l'établisse. Il est daté à la
première personne : « Hæc in Arce Apollinari dum cum Iohanne Vigintimillio
Hieracensium Marchione degeremus olim mense Augusto 1530
speculabamur. »\footnote{« Nous méditions ces choses jadis, au mois d'août 1530,
alors que nous séjournions dans la citadelle d'Apollon (\emph{in Arce
Apollinari}) avec Giovanni Ventimiglia, marquis de Geraci. » Le lieu n'est pas
identifié ici. Après « omnia », un mot court biffé n'est pas lu ; « nram » abrège
« nostram » (vue 47, droite, numéro à l'encre 46).} Suit l'annonce déjà citée :
« Nous joindrons maintenant à ceci quelques opuscules des anciens qui regardent la
même théorie sphérique. » Cette date est celle que le texte se donne ; elle ne
dit rien, et cette lecture ne dit rien, de l'antériorité de ces règles sur celles
d'autres auteurs.

\section*{La dédicace à Octavio Spinola}

\pagerange{47}{47}
Sur la même page, une lettre « D. Octavio Spinulæ Quæstori ærario Siciliæ,
Regioque Consiliario, Maurolycus. Salutem. », au trésorier de Sicile et
conseiller du roi. Il y a trois ans, écrit Maurolico, qu'on a commencé d'imprimer
à Messine le volume de ses \emph{Sphériques} ; les guerres et leurs soupçons ne
laissaient aucune place aux spéculations, qui demandent la retraite et le loisir ;
l'affection de Spinola pour les lettres a fait reprendre et achever l'ouvrage, et
« notre sphère » lui doit autant qu'Atlas dut à Hercule quand celui-ci porta le ciel
sur ses épaules. Les \emph{Sphériques} de Théodose, de Ménélas et de Maurolico
achevés, il y joindra ce qui touche la sphère mobile et le premier mouvement,
« premiers rudiments de la science des astres ».\footnote{« dabitur ut spero
quandoque ocium » : « quandoque » est écrit en abrégé, « qdoq » avec un signe
final (vue 47, droite).} Un « Hexastichon » au même, et la date : « Messanæ. Kal.
Junij M.D.LVII. », Messine, 1\textsuperscript{er} juin 1557.\footnote{« Ces
monuments longtemps troublés par diverses tempêtes courent, sous ta conduite, vers
leur but. Ainsi jadis, pour qu'une si grande machine ne pérît pas, Atlas soutint
les astres avec l'aide d'Hercule. Si ma page peut quelque chose, Octave, aucun
temps n'effacera la louange de ton mérite. » Les vers vont à la ligne un par un,
les pairs en retrait ; « Disturbata » et « procellis » sont surchargés, le « D »
de la date repassé. Les trois dates du volume — la lettre à Charles Quint de
juillet 1556, celle-ci du 1\textsuperscript{er} juin 1557, celle au vice-roi
d'août 1558, avec l'impression d'août 1558 — sont données comme le feuillet les
porte ; cette lecture ne les raccorde pas.}

\end{document}
